\documentclass[11pt, a4paper]{article}

\def\ncourse {Algebraic Geometry}
\def\ntype {problems}

\makeatletter

% packages
\usepackage{amssymb,amsfonts,amsmath,calc,tikz,pgfplots,geometry,mathtools}
\usepackage{color}   % May be necessary if you want to color links
\usepackage[svgnames]{xcolor}
\usepackage[hidelinks]{hyperref}
\usepackage{forest}
\usepackage{commath} % ew
\usepackage{esvect} % \vv makes vectors
\usepackage{centernot} % for the comparison
\usepackage{esdiff}
\usepackage{esint}
\usepackage{amsthm}
\usepackage{fancyhdr}
\usepackage{bm}
\usepackage{witharrows}
\usepackage{bookmark}
\usepackage{tikz-cd}
\usepackage{quiver}
\usepackage{bbm}
\usepackage{textcomp}
\usepackage{float}
\usepackage{gensymb}
\usepackage{cleveref}
\usepackage{environ}
\usepackage{comment}
\usepackage{cancel}
\usepackage{xparse}

% Inevitable cs
\usepackage{listings}
\usepackage[]{algorithm2e}

% tikz libraries
\usetikzlibrary{positioning}
\usetikzlibrary{matrix}
\usetikzlibrary{arrows}
\usetikzlibrary{bending}
\usetikzlibrary{arrows.meta}
\usetikzlibrary{decorations.markings}
\usetikzlibrary{decorations.pathreplacing}
\usetikzlibrary{decorations.markings}
\usetikzlibrary{patterns}
\usetikzlibrary{backgrounds}

% Page style setup
\pagestyle{fancy}
\fancyhfoffset{0pt}
\renewcommand{\headwidth}{\textwidth}
\geometry{margin=1in}
\pgfplotsset{compat=1.18}
\setlength{\headheight}{14.6pt}
\addtolength{\topmargin}{-1.6pt}
\hypersetup{
    colorlinks=false,
    linktoc=section,
    linkcolor=black,
}

%% maketitle setup
\ifx \nauthor\undefined
  \def\nauthor{Yeheli Fomberg}
\else
\fi

\ifx \nemail\undefined
  \def\nemail{\href{mailto:yp@campus.technion.ac.il}
  {\texttt{<yp@campus.technion.ac.il>}}}
\else
\fi

\ifx \ncoursehead \undefined
  \def\ncoursehead{\ncourse}
\fi

\lhead{\emph{\nouppercase{\leftmark}}}
\ifx \nextra \undefined
  \rhead{
    \ifnum\thepage=1
    \else
      \ncoursehead
    \fi}
\else
  \rhead{
    \ifnum\thepage=1
    \else
      \ncoursehead \ (\nextra)
    \fi}
\fi

\def\notes{notes}
\def\problems{problems}

\ifx \ntype \undefined
  \def\ntype{notes}
\else
\fi

\ifx \ntype \notes
  \let\@real@maketitle\maketitle
  \renewcommand{\maketitle}{\@real@maketitle\begin{center}
  \begin{minipage}[c]{0.9\textwidth}\centering\footnotesize
  These notes are not endorsed by the lecturers.
  I have revised them outside lectures to incorporate supplementary
  explanations, clarifications, and material for fun.
  While I have strived for accuracy, any errors or misinterpretations 
  are most likely mine.
  \end{minipage}\end{center}}
\else
  \ifx \ntype \problems
    \let\@real@maketitle\maketitle
    \renewcommand{\maketitle}{\@real@maketitle\begin{center}
    \begin{minipage}[c]{0.9\textwidth}\centering\footnotesize
    These problems are mostly based on the problem sets I was given
    when taking the course.
    
    While I have strived for accuracy and completeness,
    some of the problems would not have been here without the help of
    friends.
    \end{minipage}\end{center}}
  \else
  \fi
\fi


% theorem environments
\theoremstyle{definition}
\newtheorem{definition}{Definition}[section]
\newtheorem{remark}{Remark}[section]
\newtheorem{example}{Example}[section]
\newtheorem{exercise}{Exercise}[section]
\newtheorem{paradox}{Paradox}[section]
\newtheorem{entry}{Entry}[section]
\newtheorem*{solution}{Solution}
\newtheorem*{question}{Question}
\newtheorem*{answer}{Answer}
\newtheorem*{problem}{Problem}
\theoremstyle{plain}
\newtheorem{theorem}{Theorem}[section]
\newtheorem{proposition}[theorem]{Proposition}
\newtheorem{lemma}[theorem]{Lemma}
\newtheorem{corollary}[theorem]{Corollary}
\newenvironment{proofof}[1]{%
    \renewcommand{\proofname}{Proof of #1}%
    \begin{proof}
}{%
    \end{proof}
}
% tikz customization
\pgfarrowsdeclarecombine{twolatex'}{twolatex'}{latex'}{latex'}{latex'}{latex'}
\tikzset{->/.style = {decoration={markings,
                                  mark=at position 0.999
                                  with {\arrow[scale=2]{latex'}}},
                      postaction={decorate}}}
\tikzset{<-/.style = {decoration={markings,
                                  mark=at position 0.001 with {\arrowreversed[scale=2]{latex'}}},
                      postaction={decorate}}}
\tikzset{-->/.style = {decoration={markings,
                                  mark=at position 0.999
                                  with {\arrow[scale=2]{latex'[sep=0pt -2.5]}}},
                      postaction={decorate}}}
\tikzset{<--/.style = {decoration={markings,
                                  mark=at position 0.001
                                  with {\arrowreversed[scale=2]{latex'[sep=0pt -2.5]}}},
                      postaction={decorate}}}
\tikzset{<->/.style = {decoration={markings,
                                  mark=at position 0.001 with {\arrowreversed[scale=2]{latex'}},
                                  mark=at position 0.999 with {\arrow[scale=2]{latex'}},},
                      postaction={decorate}}}
\tikzset{->-/.style = {decoration={markings,
    mark=at position 0.50+0.225/#1 with {\arrow[scale=2]{latex'}}},
                      postaction={decorate}}}
\tikzset{-<-/.style = {decoration={markings,
                                   mark=at position 0.50-0.225/#1 with {\arrowreversed[scale=2]{latex'}}},
                       postaction={decorate}}}
\tikzset{->>/.style = {decoration={markings,
                                  mark=at position 1 with {\arrow[scale=2]{latex'}}},
                      postaction={decorate}}}
\tikzset{<<-/.style = {decoration={markings,
                                  mark=at position 0 with {\arrowreversed[scale=2]{twolatex'}}},
                      postaction={decorate}}}
\tikzset{<<->>/.style = {decoration={markings,
                                   mark=at position 0 with {\arrowreversed[scale=2]{twolatex'}},
                                   mark=at position 1 with {\arrow[scale=2]{twolatex'}}},
                       postaction={decorate}}}
\tikzset{->>-/.style = {decoration={markings,
                                   mark=at position 0.50+0.450/#1 with {\arrow[scale=2]{twolatex'}}},
                       postaction={decorate}}}
\tikzset{-<<-/.style = {decoration={markings,
                                   mark=at position 0.50-0.450/#1 with {\arrowreversed[scale=2]{twolatex'}}},
                       postaction={decorate}}}

\pgfarrowsdeclare{biggertip}{biggertip}{%
  \setlength{\arrowsize}{1pt}
  \addtolength{\arrowsize}{.1\pgflinewidth}
  \pgfarrowsrightextend{0}
  \pgfarrowsleftextend{-5\arrowsize}
}{%
  \setlength{\arrowsize}{1pt}
  \addtolength{\arrowsize}{.1\pgflinewidth}
  \pgfpathmoveto{\pgfpoint{-5\arrowsize}{4\arrowsize}}
  \pgfpathlineto{\pgfpointorigin}
  \pgfpathlineto{\pgfpoint{-5\arrowsize}{-4\arrowsize}}
  \pgfusepathqstroke
}
\tikzset{
	EdgeStyle/.style = {>=biggertip}
}

\tikzset{circ/.style = {fill, circle, inner sep = 0, minimum size = 3}}
\tikzset{scirc/.style = {fill, circle, inner sep = 0, minimum size = 1.5}}
\tikzset{mstate/.style={circle, draw, black, text=black, minimum width=0.7cm}}

\tikzset{eqpic/.style={baseline={([yshift=-.5ex]current bounding box.center)}}}

\definecolor{mblue}{rgb}{0.2, 0.3, 0.8}
\definecolor{morange}{rgb}{1, 0.5, 0}
\definecolor{mgreen}{rgb}{0, 0.4, 0.2}
\definecolor{mred}{rgb}{0.5, 0, 0}

% topology
\DeclareMathOperator{\dfr}{dfr} % deformation retract
\newcommand{\Cells}{\text{Cells}}
\newcommand{\RP}{\mathbb{RP}}
\newcommand{\CP}{\mathbb{CP}}

% order theory
\newcommand{\join}{\vee}
\newcommand{\meet}{\wedge}

% algebraic geometry
\DeclareMathOperator{\Rad}{Rad} % Radical of modules
\DeclareMathOperator{\Spec}{Spec}
% \renewcommand{\div}{\mathrm{div}} % divisors
\DeclareMathOperator{\Mer}{Mer} % Meromorphic functions
\DeclareMathOperator{\Exc}{Exc} % Exceptional divisor
\DeclareMathOperator{\Td}{Td} % Todd classes

% algebra

\DeclareMathOperator{\trdeg}{tr.deg.}
\DeclareMathOperator{\odd}{odd}
\DeclareMathOperator{\even}{even}
\DeclareMathOperator{\lcm}{lcm}
\DeclareMathOperator{\ord}{ord}
\DeclareMathOperator{\codim}{codim}
\DeclareMathOperator{\Ann}{Ann}
\DeclareMathOperator{\Ext}{Ext}
\DeclareMathOperator{\Sym}{Sym}
\DeclareMathOperator{\Out}{Out}
\DeclareMathOperator{\Aut}{Aut}
\DeclareMathOperator{\End}{End}
\DeclareMathOperator{\Inn}{Inn}
\DeclareMathOperator{\Mat}{Mat}
\DeclareMathOperator{\Fix}{Fix}
\DeclareMathOperator{\std}{std}
\DeclareMathOperator{\sgn}{sgn}
\DeclareMathOperator{\adj}{adj}
\DeclareMathOperator{\id}{id}
\DeclareMathOperator{\op}{op}
\DeclareMathOperator{\tr}{tr}
\DeclareMathOperator{\diag}{diag}
\DeclareMathOperator{\rank}{rank}
\DeclareMathOperator{\rk}{rk}
\DeclareMathOperator{\coker}{coker}
\DeclareMathOperator{\Mult}{Mult} % multiplier algebra
\DeclareMathOperator{\Frac}{Frac} % fraction field
\DeclareMathOperator{\Rat}{Rat}
\DeclareMathOperator{\Gal}{Gal}
\newcommand{\ioplus}{\overset{\text{int}}{\oplus}} % interior direct sum
\newcommand{\GG}{\mathcal G} % Galois correspondence
\newcommand{\FF}{\mathcal F} % Galois correspondence
\newcommand{\cupp}{\smile} % cup product
\newcommand{\capp}{\frown} % cap product
\newcommand{\actson}{\curvearrowright}
\newcommand{\bra}{\left\langle}
\newcommand{\ket}{\right\rangle}
\newcommand{\idealin}{\triangleleft\,}
\newcommand{\normalin}{\triangleleft\,}
\newcommand{\ip}[2]{\left\langle #1, #2 \right\rangle}
\newcommand{\bigslant}[2]
{{\raisebox{.2em}{$#1$}\left/\raisebox{-.2em}{$#2$}\right.}}
\newcommand{\properideal}{%
  \mathrel{\ooalign{$\lneq$\cr\raise.22ex\hbox{$\lhd$}\cr}}}

% categories
\newcommand{\cat}[1]{\mathbf{#1}}
\newcommand{\Ch}{\cat{Ch}} % Chain complexes   
\newcommand{\Set}{\cat{Set}}
\newcommand{\Grp}{\cat{Grp}}
\newcommand{\Ab}{\cat{Ab}}
\newcommand{\Ring}{\cat{Ring}}
\newcommand{\CRing}{\cat{CRing}}
\newcommand{\Top}{\cat{Top}}
\newcommand{\Vect}{\cat{Vect}}
\newcommand{\cmod}{\cat{mod}}

%% matrix groups
\newcommand{\GL}{\mathrm{GL}}
\newcommand{\Or}{\mathrm{O}}
\newcommand{\PGL}{\mathrm{PGL}}
\newcommand{\PSL}{\mathrm{PSL}}
\newcommand{\PSO}{\mathrm{PSO}}
\newcommand{\PSU}{\mathrm{PSU}}
\newcommand{\SL}{\mathrm{SL}}
\newcommand{\msl}{\mathrm{sl}}
\newcommand{\SO}{\mathrm{SO}}
\newcommand{\Spin}{\mathrm{Spin}}
\newcommand{\Sp}{\mathrm{Sp}}
\newcommand{\SU}{\mathrm{SU}}
\newcommand{\U}{\mathrm{U}}
\let\char\relax
\newcommand{\char}{\text{char}}

% analysis
\renewcommand{\Re}{\mathrm{Re}}
\renewcommand{\Im}{\mathrm{Im}}
\NewDocumentCommand{\dx}{o}{
  \IfNoValueTF{#1}
    {\dif x}                  % if no optional argument
    {\dif^{#1} \!x}         % if optional argument
}
\NewDocumentCommand{\dy}{o}{
  \IfNoValueTF{#1}
    {\dif y}                  % if no optional argument
    {\dif^{#1} \!y}         % if optional argument
}
\newcommand{\dt}{\dif t}
\newcommand{\du}{\dif u}
\newcommand{\dv}{\dif v}
\newcommand{\dz}{\dif z}
\newcommand{\ds}{\dif s}
\newcommand{\dw}{\dif w}
\newcommand{\dr}{\dif r}
\newcommand{\dm}{\dif m}
\newcommand{\df}{\dif f}
\newcommand{\dV}{\dif V}
\newcommand{\dS}{\dif S}
\newcommand{\dA}{\dif \mathrm{A}}
\newcommand{\dmu}{\dif \mu}
\newcommand{\dnu}{\dif \nu}
\newcommand{\dtheta}{\dif \theta}
\newcommand{\domega}{\dif \omega}
\newcommand{\dsigma}{\dif \sigma}
\newcommand{\dtau}{\dif \tau}
\newcommand{\dvarphi}{\dif \varphi}
\renewcommand{\dif}{\mathrm{d}}
\renewcommand{\Dif}{\mathrm{D}}
\renewcommand{\triangle}{\Delta}
\newcommand{\ptriangle}{\partial \triangle}
\DeclareMathOperator{\im}{im}
\DeclareMathOperator{\cis}{cis}
\DeclareMathOperator{\rot}{rot}
\DeclareMathOperator{\vspan}{span}
\DeclareMathOperator{\grad}{grad}
\DeclareMathOperator{\curl}{curl}
\DeclareMathOperator{\esssup}{esssup}
\newcommand{\hodge}{{\mathnormal{\star}}}
\DeclareMathOperator{\range}{range}
\DeclareMathOperator{\Hom}{Hom}
\DeclareMathOperator{\Int}{Int}
\DeclareMathOperator{\Conv}{Conv} % convex hull
\newcommand{\ind}{\mathbbm{1}}
\DeclareMathOperator{\Res}{Res}
\DeclareMathOperator{\graph}{graph}
\DeclareMathOperator{\diam}{diam}
\DeclareMathOperator{\supp}{supp}
\DeclareMathOperator{\Vol}{Vol} % Volume
\DeclareMathOperator{\Area}{Area} % Area
\DeclareMathOperator{\area}{area}
\DeclareMathOperator{\Ind}{Ind} % Index
\newcommand{\length}{\mathrm{length}}

% logic
\DeclareMathOperator{\MOD}{MOD}
\DeclareMathOperator{\Theory}{Theory}
\newcommand{\notimplies}{%
  \mathrel{{\ooalign{\hidewidth$\not\phantom{=}$\hidewidth\cr$\implies$}}}}

% nice
\newcommand{\half}{\frac{1}{2}}
\newcommand{\pair}{\del}
\newcommand{\taking}[1]{\xrightarrow{#1}}
\newcommand{\otaking}[1]{\xleftarrow{#1}}
\newcommand{\inv}{^{-1}}
\newcommand{\ot}{\leftarrow}
\newcommand{\ninfty}{-\infty}
\newcommand{\pinfty}{+\infty}
\newcommand{\floor}[1]{\left\lfloor #1 \right\rfloor}
\newcommand{\ceil}[1]{\left\lceil #1 \right\rceil}
\newcommand{\viff}{\Big\Updownarrow} % vertical iff
\newcommand{\bmid}{\,\middle\vert\,} % big mid

% probability
\newcommand{\Prob}{\mathbf{P}}
\renewcommand{\vec}[1]{\boldsymbol{\mathbf{#1}}}
\DeclareMathOperator{\Bin}{Bin}
\DeclareMathOperator{\Geo}{Geo}
\DeclareMathOperator{\Poi}{Poi}
\DeclareMathOperator{\Exp}{Exp}
\DeclareMathOperator{\Var}{Var} % Variance
\DeclareMathOperator{\Cov}{Cov}

% special letters
\newcommand{\N}{\mathbb{N}}
\newcommand{\Z}{\mathbb{Z}}
\newcommand{\Q}{\mathbb{Q}}
\newcommand{\R}{\mathbb{R}}
\newcommand{\C}{\mathbb{C}}
\newcommand{\F}{\mathbb{F}}
\newcommand{\E}{\mathbf{E}}
\newcommand{\D}{\mathbb{D}}
\renewcommand{\H}{\mathbb{H}}
\newcommand{\ps}{\mathcal{P}}
\newcommand{\M}{\mathcal{M}}
\renewcommand{\L}{\mathcal{L}}
\renewcommand{\O}{\mathcal{O}}
\renewcommand{\U}{\mathcal{U}}
\newcommand{\Omicron}{O}
\newcommand{\powerset}{\mathcal{P}}

% text
\newcommand{\st}{\text{ s.t. }}
\newcommand{\tand}{\quad \text{and} \quad}
\newcommand{\tor}{\quad \text{or} \quad}
\newcommand{\stand}{\text{ and }}
\newcommand{\stor}{\text{ or }}
\renewcommand{\tt}[1]{\textnormal{\textbf{(#1).}}} %tt=theorem title GET RID OF

% title format
\title{\textbf{\ncourse}}

\ifx \ntype \notes
  \author{Notes taken by \nauthor\, \nemail \\ Based on lectures by \nlecturer}
  \date{\nterm\ \nyear}
\else
  \ifx \ntype \problems
    \author{These problems were solved by \nauthor \\ \nemail}
  \else
  \fi
\fi

\makeatother

\renewcommand{\P}{\mathbb{P}}
\renewcommand{\div}{\mathrm{div}} % divisors

\begin{document}
\tableofcontents
\newpage

\section{Algebraic subsets of \texorpdfstring{$\C^n$}{C}}

\begin{exercise}
  Prove that $\sqrt{\sqrt{I}} = \sqrt{I}$.
\end{exercise}
\begin{solution}
  Recall the radical of an ideal $I \normalin R$
  is given by:
  \[
    \sqrt{I} :=
    \set{r \in R \mid r^n \in I \text{ for some } n \in \Z^{+}}.
  \]
  We will prove this by double-sided containment.
  Let $a \in \sqrt{I}$.
  Then for $1 \in \Z^+$ we have $a^1 \in \sqrt{I}$,
  and therefore by definition $a \in \sqrt{\sqrt{I}}$.

  Next let $a \in \sqrt{\sqrt{I}}$.
  By definition there exists $n_1 \in \Z^+$ such that $a^{n_1} \in \sqrt{I}$,
  which again by definition means that there is some $n_2 \in \Z^+$
  such that $\del{a^{n_1}}^{n_2} \in I$.
  We get that $a^{n_1 n_2} \in I$, and thus $a \in \sqrt{I}$
  which completes the proof.
\end{solution}

\begin{exercise}
  Prove that $Z(I(V)) = \overline V$ (the closure in the Zariski topology).
\end{exercise}
\begin{solution}
  Recall that the closure of a set $V$ is the intersection of all closed
  sets that contain in.
  Equivalently, it is the smallest closed set that contains $V$.

  We know that $V \subset Z(I(V))$ from the theorem that $(I,Z)$
  is a Galois connection (this also follows directly from the definition).
  It is left to show that $Z(I(V))$ is the smallest closed subset containing
  $V$.

  Let $V \subseteq C = Z(J)$ be a Zariski closed subset that contains $V$.
  By definition of $Z$, it follows that all polynomials in $J$
  vanish on $V$, and therefore $J \subseteq I(V)$.
  By the Galois connection we know that $Z$ inverses relations,
  and thus $Z(I(V)) \subseteq Z(J) = C$.
  This implies that $Z(I(V))$ is the smallest closed subset containing $V$,
  which completes the proof.
\end{solution}

\begin{exercise}
  Find counterexamples to Theorems 2.5, 2.7 and Corollary 2.6 for $\C$
  replaced by $\R$ in the notes.
\end{exercise}
\begin{solution}
  For theorem 2.5, one counterexample is the ideal $(x^2 + 1) \idealin \R[x]$.
  Since $x^2 + 1$ is irreducible over $\R$, the ideal $(x^2 + 1)$ is maximal.
  However by its irreducibility, it is clearly not a point ideal.

  For corollary 2.6, we can use the same counterexample $I := (x^2 + 1)$
  since clearly $Z(I) = \emptyset$ (over $\R$) and $I \neq (1)$.

  For theorem 2.7, we can use the same counterexample again $I = (x^2 + 1)$.
  As we have seen $Z(I) = \emptyset$, and thus $I(Z(I)) = \R[x]$.
  However, $\sqrt{I} \neq \R[x]$, since for example, $1 \notin \sqrt{I}$
  (otherwise $1 \in I$ which is not the case).
\end{solution}

\newpage

\section{Irreducible components}

\begin{exercise}
  Prove that for a proper ideal $I \subset \C[x_1,\dots,x_n]$ the set
  $X := Z(I) \subseteq \C^n$ is irreducible if and only if $\sqrt{I}$ is a prime
  ideal.
\end{exercise}
\begin{solution}
  By the regular strength nullstellensatz theorem we get that
  $I(Z(I)) = \sqrt{I}$, so we can replace $\sqrt{I}$ in the theorem
  with $I(Z(I))$.

  First assume that $I(Z(I))$ is not a prime ideal.
  This means that exists $f,g$ different than $1$ that don't vanish on $Z(I)$,
  such that $fg$ vanishes on $Z(I)$.
  Since $fg$ vanishes on $Z(I)$ we get that $Z(I) \subset Z(fg) = Z(f) \cup Z(g)$,
  so we get that
  \[
    Z(I) = \del{X \cap Z(f)} \cup \del{X \cap Z(g)}.
  \]
  It is clear that $\emptyset \neq X \cap Z(f)$,
  since otherwise either $g$ vanishs on $X$ which is a contradiction,
  and also $X \cap Z(f) \subsetneq X$,
  since $f$ vanishes on $X \cap Z(f)$ but not on $X$.
  Similarly we get that $\emptyset \neq X \cap Z(g) \subsetneq X$,
  thus $X := Z(I)$ is reducible.

  Next assume that $Z(I)$ is reducible.
  Suppose that its decomposition is $Z(I) = Z(J_1) \cup Z(J_2)$.
  We get that
  \[
    I(Z(I)) =
    I\del{Z(J_1) \cup Z(J_2)} =
    I\del{Z(J_1 \cap J_2)}.
  \]
  By the Galois connection notice that we have
  \[
    I(Z(J_1)) \supsetneq I(Z(I)) \tand
    I(Z(J_2)) \supsetneq I(Z(I)).
  \]
  Choose $f \in I(Z(J_1)) \setminus I(Z(I))$ and
  $g \in I(Z(J_1)) \setminus I(Z(I))$.
  We now get that
  \[
    fg \in I\del{Z(J_1 \cap J_2)} = I(Z(I)),
  \]
  which implies that $I(Z(I))$ is not prime, and completes the proof.
\end{solution}

\begin{exercise}
  Prove that $Z(I)$ is irreducible if and only if any nonempty open subset of
  $Z(I)$ is dense (with respect to the induced topology).
\end{exercise}
\begin{solution}
  First let us prove that any open nonempty subset of an irreducible closed set
  $Z(I)$ is dense.
  Let $U$ be an open nonempty set such that $\overline U \neq Z(I)$.
  Then immediately we have $Z(I) = U^c \cup \overline U$ in contradiction to $Z(I)$
  being irreducible.
  
  Next, suppose that any nonempty subset of $Z(I)$ is dense.
  Assume, by way of contradiction, that $Z(I)$ is reducible with
  decomposition $Z(I) = Z(J_1) \cup Z(J_2)$.
  Denote the open (in the induced topology) set $D := Z(I) \setminus Z(J_1)$.
  We have that:
  \[
    D = Z(I) \setminus Z(J_1) \subseteq Z(J_2).
  \]
  By the assumption of the exercise $\overline D = Z(I)$,
  however $D \subset B$ which is closed so
  $\overline D \subset B \subsetneq I(Z)$.
  This contradiction completes the proof.
\end{solution}

\begin{exercise}
  Verify that the transition functions for the cover of $\CP^n$
  by $U_i$ are indeed holomorphic.
\end{exercise}
\begin{solution}
  The charts $\varphi_i \colon U_i \to \C^n$ are given by:
  \[
    \varphi_i([x_0 : \cdots : x_n]) =
    \del{\frac{x_0}{x_0},\dots,\frac{x_{i-1}}{x_i},
      \frac{x_{i+1}}{x_i},\dots,\frac{x_n}{x_i}}.
  \]
  We must check that the transition functions $\varphi_j \circ \varphi_i^{-1}$
  are holomorphic on the intersection $\varphi_i(U_i \cap U_j)$.
  Let $z = (z_0,\dots,z_{i-1},z_{i+1},\dots,z_n) \in \varphi_i(U_i \cap U_j) \subset
  \C^n$.
  Then the preimage of it by $\varphi_i$ is
  the point $[z_0,\dots,z_{i-1},1,z_{i+1},\dots,z_n]$.
  However it must also point in $U_j$ so we may divide by $z_j$ and get:
  \[
    \varphi_j \circ \varphi_i^{-1} (z) =
    \del{\frac{z_0}{z_j},\dots,\frac{z_{i-1}}{z_j},\frac{1}{z_j},
    \frac{z_{i+1}}{z_j},\dots,\frac{z_n}{z_j}}.
  \]
  Since each component is a rational function with a nonvanishing
  denominator (in the domain $\set{x_j \neq 0}$) we get that
  the transition maps are holomorphic.
\end{solution}

%% MISSING HILBERT POLYNIALS AND BEZOUTS

\newpage

\section{Line bundles and vector bundles}

\begin{exercise}
  Verify that $(6.2)$ defines a section of $\mu^*L$.
\end{exercise}
\begin{solution}
  We need to verify that $\mu^*s \colon X \to \mu^*L$ defined by
  \[
    \mu^*s \colon x \mapsto (x,s(\mu(x))) \in X \times L
  \]
  defines a section of $\mu^*L$.
  A section is a map $\mu^*s \colon X \to \mu^*L$,
  and therefore we need to verify that $\mu^*s(x) \in \mu^*L \to X$.
  Thus we need to check that:
  \[
    (x,s(\mu(x))) \in
    \set{(x,l) \in X \times L \colon \mu(x) = \pi(l)}.
  \]
  Indeed, since $s$ is a section on $L \to Y$ we get that
  $\pi(s(\mu(x))) = \id_Y(\mu(x)) = \mu(x)$,
  so the section is well defined.
  Clearly it is holomorphic as composition of holomorphic functions.
  It is left to check that $\pi^* \circ s = \id_X$,
  where $\pi^*$ is the proejction map $\pi^* \colon \mu^*L \to X$
  defined by $\pi^*(x,l) = x$.
  However, this is easy to verify since
  \[
    \pi^*(\mu^*s(x)) =
    \pi^*(x,s(\mu(x))) = x,
  \]
  so clearly $\pi^* \circ \mu^*s = \id_X$.
\end{solution}

\begin{exercise}
  Show that sections of a vector bundle $W \to X$ have a
  natural structure of a vector space over $\C$.
\end{exercise}
\begin{solution}
  Let $p \colon W \to X$ be a vector bundle as defined in class.
  We need to show that the space of global sections of $L$ $\Gamma(X,W)$
  has the natural structure of a complex vector space.
  We define addition and scalar multiplication pointwise.
  That is, for any $s,t \in \Gamma(X,W)$ and $\lambda \in \C$ we have:
  \[
    (s + t)(x) := s(x) + t(x) \in p^{-1}(x) \tand
    (\lambda s)(x) := \lambda\del{s(x)} \in p^{-1}(x),
  \]
  where the inclusions in the fiber $p^{-1}(x)$ follows from the fact
  it is a complex vector space.
  Since the projection map $p$ is linear we also have
  $p \circ (s + t) = \id_X$ and $p \circ (\lambda s)_X$
  so indeed $(s+t),(\lambda s) \in \Gamma(X,W)$
  (this means that addition and multiplication
  by scalar are well defined on $\Gamma(X,W)$).

  Commutativity follows from commutativity of pointwise function addition,
  and similarly associativity, distributivity, and the rest of the axioms
  of a vector space still hold.

  We will in particular mention that the zero vector
  is the zero section $0 \in \Gamma(X,W)$
  defined by $0(x) = 0_{p^{-1}(x)}$ (the section that sends each point
  in $X$ to the zero of its corresponding vector space),
  since $(s + 0)(x) = s(x) + 0(x) = s(x)$.

  Similarly, the additive inverse of $s \in \Gamma(X,W)$ is $(-s)$,
  given by $(-s)(x) = -s(x)$.
  It is a well defined section since $s(x) \in p^{-1}(x)$ is a vector in a
  complex vector space and has a well defined additive inverse.
\end{solution}

\begin{exercise}
  Verify that sections of the line bundle $\C \times X \to X$ are in
  natural bijection to holomorphic functions on $X$.
\end{exercise}
\begin{solution}
  Every section $s \in \Gamma(X,\C \times X)$ can be written
  as $s(x) = (f(x),x)$ where $f$ is some map from $x$ to the first
  coordinate of $s$.
  Since as we wrote $p \circ s = \id_X$, it follows that a section is
  \textit{completely determined} by the map $f \colon X \to \C$.
  Conversely, any such map \textit{uniquely} defines a section by $s(x) = (f(x),x)$.
  This shows that there is a bijection between sections of the line bundle
  and some maps $f$ on $X$.

  In the case of a section $s \in \Gamma(X,\C \times X)$,
  since a section is required to be holomorphic (by definition),
  the map it gives rise to $f \colon X \to \C$ must also be holomorphic.
  Conversely, a holomorphic map $f \colon X \to \C$ gives rise
  to a section $s(x) = (f(x),x)$.
  Using this correspondence and the uniqueness discussed in the previous
  paragraph, we conclude that sections of the line bundle $\C \times X \to X$
  are in natural bijection to holomorphic functions on $X$
  which is given by
  the corrspondence
  $\Gamma(X,\C \times X) \ni s \longleftrightarrow f \in \mathrm{Hol}(X)$.
\end{solution}

\newpage

\section{More line bundles and maps to projective spaces}

\begin{exercise}
  Verify that transition functions for the identification of $(7.1)$
  (but for all $U_i$) are holomorphic and rational on $U_i \cap U_j$
  and thus give $\O(1)$ the structure of a line bundle over $\CP^n$
  in both holomorphic or algebraic settings.
\end{exercise}
\begin{solution}
  Recall that the bijection $\pi^{-1} U_0 \cong \C \times U_0$
  is given by:
  \[
    \del{[x_0 : \cdots : x_n],\varphi} \mapsto
    \del{\varphi\del{1,\frac{x_1}{x_0},\dots,\frac{x_n}{x_0}},[x_0 : \cdots : x_n]}.
  \]
  Fix some arbitrary $i$,
  and let the map $v^{(i)} \colon U_i \to \CP^{n+1}$
  be the map that normalizes a point in projective coordinates
  relative to that $i$th coordinate:
  \[
    v^{(i)}([x_0 : \cdots : x_n]) =
    \del{\frac{x_0}{x_i},\dots,\frac{x_{i-1}}{x_i},1,\frac{x_{i+1}}{x_i},\dots,\frac{x_n}{x_i}}
    \in \C^{n+1}.
  \]
  Define the $i$th trivialization map
  $\Psi_i \colon \pi^{-1}(U_i) \to \C \times U_i$ to be
  as shown in the notes and beginning of this solution:
  \[
    \Psi_i([x_0 : \cdots : x_n],\varphi) =
    \del{\varphi\del{v^{(i)}([x_0 : \cdots : x_n])},[x_0 : \cdots : x_n]}.
  \]
  The transition maps are the maps of the form $\Psi_i \circ \Psi_j^{-1}$.
  Note that since $\varphi$ is a linear functional (an element in the dual
  space of the line $l_p$),
  we have the relation:
  \[
    v^{(i)}\left([x_0 : \cdots : x_n]\right) =
    \frac{x_j}{x_i} v^{(j)}\left([x_0 : \cdots : x_n]\right),
  \]
  because multiplying $v^{(j)}([x_0 : \cdots : x_n])$
  by $\frac{x_j}{x_i}$ makes the $i$th coordinate $1$ and the other coordinates
  agree with those of $v^{(i)}([x_0 : \cdots : x_n])$.
  Therefore the transition maps
  $\Psi_i \circ \Psi_j^{-1} \colon \C \times (U_i \cap U_j) \to \C \times (U_i \cap U_j)$
  are simply:
  \[
    \Psi_i \circ \Psi_j^{-1} \colon
    \left(a,[x_0 : \cdots : x_n]\right) \mapsto
    \left(\frac{x_j}{x_i} a,[x_0 : \cdots : x_n]\right).
  \]
  The function $\frac{x_j}{x_i}$ is nowhere zero on the intersection
  $U_i \cap U_j$, and thus holomorphic.
  It is also a rational function since it is a ratio of homogeneous
  polynomials of the same degree.
\end{solution}

\begin{exercise}[Bonus exercise]
  Prove that every holomorphic section of $\O(1) \to \CP^n$
  is a linear combination of homogeneous coordinates.
  Hint: for any such
  section define the map $f \colon (\C^{n+1} \setminus \set{0}) \to \C$
  and show $f$ is holomorphic.
  Use Hartogs Lemma to extend $f$ to a holomorphic function on $\C^{n+1}$ and
  use Taylor expansion and $f(\lambda z) = \lambda z$.
  The same argument shows that
  holomorphic sections of $\O(k) \to \CP^n$ are zero for $k < 0$
  are are homogeneous polynomials of degree $k$ for $k \geq 0$.
\end{exercise}
\begin{solution}
\end{solution}

\begin{exercise}
  Prove that the maps $(7.2)$ and $(7.3)$ are embeddings of
  complex manifolds.
\end{exercise}
\begin{solution}
  The map given in $(7.2)$ is the map $f \colon \CP^1 \to \CP^2$ defined:
  \[
    [y_0 : y_1] \mapsto [y_0^2 : y_0 y_1 : y_1^2].
  \]
  Holomorphicity of this map is given by the homogeneous
  quadratic polynomials in the homogeneous coordinates.
  It is also injective by matter of pure set theoretic calculation.

  We see that on the affine chart $U_0 = \set{y_0 \neq 0}$ with
  the coordinate $t = y_1/y_0$ the local representative of $f$
  is
  \[
    t \mapsto (1,t,t^2).
  \]
  Its derivative is $(0,1,2t)$, which is never the zero tangent vector.
  This implies that the differential is injective at every point so $f$
  is an immersion.
  Clearly this works for the other chart as well.
  Since $\CP^1$ is compact $f$ is automatically proper.
  By a theorem, if $f$ is injective, proper, and an immersion (full rank
  differential), it is an embedding.
  
  The map given in $(7.3)$ is the map $f \colon \CP^1 \times \CP^1 \to \CP^3$
  defined:
  \[
    ([y_0 : y_1],[z_0 : z_1]) \mapsto [y_0 z_0 : y_0 z_1 : y_1 z_0 : y_1 z_1].
  \]
  It is holomorphic by the same argument
  from earlier and it is also injective by matter of pure set
  theoretic calculation.

  Work on the affine chart where $y_0 \neq 0$ and $z_0 \neq 0$.
  Put local coordinates:
  \[
    s = \frac{y_{1}}{y_{0}} \tand
    t = \frac{z_{1}}{z_{0}}.
  \]
  Then, in the affine chart of $\CP^3$ with the first coordinate scaled
  to $1$ we get that the map is
  \[
    (s,t) \mapsto (1,t,s,st).
  \]
  The Jacobian of this map is:
  \[
  \begin{pmatrix}
    0 & 0 \\
    0 & 1 \\
    1 & 0 \\
    t & s
  \end{pmatrix}
  \]
  which has degree $2$ so the differential is injective at every point
  of this chart.
  Clearly this works for the other charts as well.
  Since $\CP^1$ is compact and product of compact spaces is compact
  we get that $\CP^1 \times \CP^1$ is compact and thus the map is also proper.
  By a theorem, if $f$ is injective, proper, and an immersion (full rank
  differential), it is an embedding.
\end{solution}

\newpage

\section{Weil and Cartier divisors and class groups}

\begin{exercise}
  Find the Weil divisor $\div f$ for the rational function on $\C^2$
  \[
    f(x_1,x_2) =
    \frac{x_1^3 x_2}{x_1^2 + x_2^3 - x_2}.
  \]
  % hint: it is simple
\end{exercise}
\begin{solution}
  We write $f = \frac{N}{D}$ with
  \[
    N(x_1,x_2) = x_1^3 x_2 \tand
    D(x_1,x_2) = x_1^2 + x_2^3 - x_2.
  \]
  The Weil divisor of $f$ over a variety $X$ is thus:
  \[
    \div f =
    \div N - \div D =
    \sum_{Y} \ord_Y(N) Y - \sum_{Y} \ord_Y(D) Y,
  \]
  where we sum over all prime divisors $Y$ of $X$,
  and $\ord_Y(f) Y$  is the order of $f$ along $Y$.
  It is clear the the pirime divisors of $N$ are exactly
  $\set{x_1 = 0}$ and $\set{x_2 = 0}$ since $x_1$ and $x_2$
  are the only irreducible factors of $N$.
  It follows that:
  \[
    \div N =
    \sum_{Y} \ord_Y(N) Y =
    3 \set{x_1 = 0} + \set{x_2 = 0}.
  \]
  Similarly, we see that $x_1^2 + x_2^3 - x_2$ is irreducible in $\C[x_1,x_2]$,
  and since the order is $1$ we get that:
  \[
    \div D =
    \sum_{Y} \ord_Y(D) Y =
    \set{x_1^2 + x_2^3 - x_2 = 0},
  \]
  and finally,
  \[
    \div f =
    3 \set{x_1 = 0} + \set{x_2 = 0} -
    \set{x_1^2 + x_2^3 - x_2 = 0}.
  \]
\end{solution}

\begin{exercise}
  Find a Cartier divisor on $\CP^1$ so that the corresponding
  Weil divisor is $5[0] - [i]$.
  % Hint: use the cover of $\CP^1$ by $U_0$ and $U_1$.
\end{exercise}
\begin{solution}
    The Cartier divsor we choose (wisely) is:
  \[
    D = \set{\del{U_0, \frac{1}{z^4 (1 - i z)}}, \del{U_1, \frac{w^5}{w - i}}},
  \]
  where $w = 1/z$.
  We see that on the intersection $U_0 \cap U_1$ we have that:
  \[
    \frac{f_0}{f_1} =
    \frac{1}{z^4 (1 - i z)} \cdot \frac{1/z - i}{(1/z)^5} =
    \frac{1}{z^4 (1 - i z)} \cdot \frac{1 - i z}{z} \cdot z^5 = 1.
  \]
  Since $f_0/f_1 = 1$ is a unit on $U_0 \cap U_1$ we get that this is indeed
  a valid Cartier divisor.
  We now have to show that its corresponding Weil divisor is the desired one.
  Indeed, since we chose $U_0$ with coordinate $z$ ($x_0 \neq 0$)
  and $U_1$ with coordinate $w$ ($x_1 \neq 0$),
  we get that for $[0] = (0 : 1)$ we have that for $\frac{w^5}{w - i}$
  (we must check this chart since $[0] \notin U_0$):
  \[
    \ord_{[0]}(\frac{w^5}{w - i}) = 5,
  \]
  and similarly for $[i] = (i : 1)$ (this time we may use any chart)
  we get:
  \[
    \ord_{[i]}(\frac{w^5}{w - i}) =
    \ord_{w = i}(\frac{w^5}{w - i}) = -1,
  \]
  and then we get that the corresponding Weil divisor is indeed
  $5[0] - [i]$, which completes the proof.
\end{solution}

\begin{exercise}
  Consider the quotient of the ring of formal power series in three variables
  \[
    R = \C[[x_1,x_2,x_3]] / (x_1 x_2 - x_3^2).
  \]
  Show that the ideal $I = (x_1, x_3) \subset R$ is not principle.
  % Hint: look at the dimension of $I / (x_1,x_2,x_3) I$.
\end{exercise}
\begin{solution}
  Let $S = \C[[x_1,x_2,x_3]]$ and so $R = S/(x_1 x_2 - x_3^2)$.
  Note that $R$ is a regular local ring with maximal
  ideal $\mathfrak m = (\bar x_1,\bar x_2,\bar x_3) R$.
  By Nakayama's lemma, if $I = (x_1,x_3)$ were principle, we must have that
  $\dim_{\C} I/\mathfrak m I = 1$.
  We will show that $\dim_{\C} I/\mathfrak m I \geq 2$ and get a contradiction.

  The ideal $I$ is generated by the images of $x_1$ and $x_3$
  in $R$ which we will denote by $\bar x_1$ and $\bar x_3$.
  We notice that $\mathfrak m I \subset \mathfrak m^2$.
  That is because $\mathfrak m I = \mathfrak m(x_1,x_3)$,
  so every element of $\mathfrak m I$ has degree at least $2$.
  Thus the images $\bar x_1$ and $\bar x_3$ cannot be zero.

  Assume, by way of contradiction, that $\bar x_1$ and $\bar x_3$
  are linearly dependent.
  That is, there exist nonzero $\alpha,\beta \in \C$ such that
  \[
    \alpha \bar x_1 + \beta \bar x_3 = 0 \quad\text{in $I/\mathfrak m I$}.
  \]
  This means that $\alpha x_1 + \beta x_3 \in \mathfrak m I$,
  however since both $x_1$ and $x_3$ are of degree $1$,
  which implies that $\alpha = \beta = 0$.
  This contradiction shows that $\dim_{\C} I/\mathfrak m I \geq 2$
  which completes the proof.
\end{solution}

\newpage

\section{Cartier divisors as meromorphic sections of line bundles}

\begin{exercise}
  Verify that Definition $9.1$ in the notes defines a Cartier divisor.
  Further verify that different choices of isomorphisms
  $\pi^{-1}(U_{\alpha}) \cong C \times U_{\alpha}$
  define the same Cartier divisor.
\end{exercise}
\begin{solution}
  It's just opening definitions but I don't have time :(
\end{solution}

\begin{exercise}
  Verify that the relation $\sim $in Definition $9.4$ in the notes is indeed an
  equivalence relation.
\end{exercise}
\begin{solution}
  For a Cartier divisor $\set{U_{\alpha},f_{\alpha}}_{\alpha \in I}$
  we define the relation $\sim$ on $\bigsqcup_{\alpha} (\C \times U_{\alpha})$
  with
  \[
    (\C \times U_{\alpha}) \ni 
    (t,x) \sim \del{\frac{f_{\beta}(x)}{f_{\alpha}(x)}t, x}
    \in (\C \times U_{\beta}).
  \]
  We will show that it is an equivalence relation.
  In what follows, we denote an element of $\C \times U_{\alpha}$
  by $(t,x)_{\alpha}$ etc.
  Note that there is a relation $(t,x)_{\alpha} \sim (s,y)_{\beta}$
  if and only if $x = y \in U_{\alpha} \cap U_{\beta}$
  and $s = \frac{f_{\beta}(x)}{f_{\alpha}(x)}t$.
  
  \noindent
  \textbf{Reflexive:}
  The goal is to show that $(t,x)_{\alpha} \sim (t,x)_{\alpha}$ for any $\alpha$.
  We just need to show that $t = \frac{f_{\alpha}(x)}{f_{\alpha}(x)}t$,
  but this is also clear so this is pretty much immediate.
  
  \noindent
  \textbf{Symmetric:}
  Given $(t,x)_{\alpha} \sim (s,y)_{\beta}$
  we show $(s,y)_{\beta} \sim (t,x)_{\alpha}$.
  The conditions for the relation are pretty much symmetric,
  and the only not exactly trivial thing is to show given that
  \[
    s = \frac{f_{\beta}(x)}{f_{\alpha}(x)}t,
  \]
  that we also have
  \[
    t = \frac{f_{\alpha}(x)}{f_{\beta}(x)}s,
  \]
  but we get this by multiplying both sides of the expression given
  by $\frac{f_{\alpha}(x)}{f_{\beta}(x)}$.
  
  \noindent
  \textbf{Transitivity:}
  It is clear that if $x = y \in U_x \cap U_y$
  and $y =  \in U_y \cap U_z$ then $x = z \in U_x \cap U_z$.
  It also follows from basic fraction multiplication that given
  \[
    s = \frac{f_{\beta}(x)}{f_{\alpha}(x)}t \tand
    r = \frac{f_{\gamma}(x)}{f_{\beta}(x)}s,
  \]
  we can conclude:
  \[
    r = \frac{f_{\gamma}(x)}{f_{\alpha}(x)}t,
  \]
  and thus transitivity follows,
  which completes the proof.
\end{solution}

\begin{exercise}
  Check that the canonical line bundle of $\CP^n$ is isomorphic
  to $\O(-n-1)$, for any $n \geq 1$.
\end{exercise}
\begin{solution}
  Let us compute the canonical line bundle $K$ of $\CP^n$.
  Rational sections of $K$ are rational differential $n$-forms,
  so we can pick one of them and compute the corresponding Cartier divisor.
  In fact, each $n$-form $f(z_1,\dots,z_n) \dz_1 \wedge \cdots \wedge \dz_n$
  will correspond to $f(z_1,\dots,z_n)$ which will give us a Cartier divisor.
  Let us start with a general form:
  \[
    \omega =
    \dif\del{\frac{x_1}{x_0}}
    \dif\del{\frac{x_2}{x_0}} \cdots
    \dif\del{\frac{x_n}{x_0}},
  \]
  which is precisely $\dz_1 \wedge \dz_2 \wedge \cdots \wedge \dz_n$,
  for the coordinates on $U_0$.
  Thus the corresponding function is $f_0 = 1$.
  We now wish to compute it on the rest of the charts.
  We see from computations similar to those we did in class
  in the $\CP^2$ case that for a chart $U_{\alpha}$
  for some natural $0 < \alpha \le n$ that:
  \begin{align*}
    \omega &=
    \dif\del{\frac{x_1}{x_{\alpha}(\frac{x_0}{x_{\alpha})^{-1}} }} \cdots
    \dif\del{\frac{x_\alpha}{x_0}}^{-1} \cdots
    \dif\del{\frac{x_n}{x_{\alpha}(\frac{x_0}{x_{\alpha})^{-1}} }} = \cdots \\ &=
    -\del{\frac{x_{0}}{x_{\alpha}}}^{-n-1} 
    \dif\del{\frac{x_0}{x_{\alpha}}} \cdots
    \widehat{\dif\del{\frac{x_{\alpha}}{x_{\alpha}}}} \cdots
    \dif\del{\frac{x_n}{x_{\alpha}}},
  \end{align*}
  which gives us $f_{\alpha} = \del{\frac{x_{0}}{x_{\alpha}}}^{-n-1}$
  (the sign does not matter since it is an invertible function).
  Thus, the corresponding Weil divisor is $(-n-1)[\set{x_0 = 0}]$,
  and thus the canonical line bundle of $\CP^n$ is isomorphic to $\O(-n-1)$
  for any $n \geq 1$.
\end{solution}

\newpage

\section{Cubic curve in \texorpdfstring{$\CP^2$}{CP}. Group law.}

\begin{exercise}
  Extend the count of the codimension of the bad locus to
  arbitrary hypersurfaces of degree $d$ in $\CP^n$ by again figuring out what it
  takes to be singular at $(1 : 0 : . . . : 0)$.
\end{exercise}
\begin{solution}
  Suppose that $\set{f = 0}$ is singular at $(1 : 0 : . . . : 0)$.
  This point lies in $U_0$ and we can dehomogenize $f = 0$ to the equation:
  \[
    \sum_{|\alpha| = d}
    c_{\alpha}
    \del{\frac{x_1}{x_0}}^{\alpha_1}
    \del{\frac{x_2}{x_0}}^{\alpha_2} \cdots
    \del{\frac{x_n}{x_0}}^{\alpha_n} = 0,
  \]
  where $\alpha$ is a multi-index with $n+1$ entries, and its first
  entry (with index $0$) is given by $d - \sum_i \alpha_i$.
  In order for this curve to pass through $(1 : 0 : . . . : 0)$ (the origin
  of $U_0$), we need to have no constant term,
  which means that $c_{(d,0,\dots,0)} = 0$.
  Moreover, if any of the linear terms is nonzero,
  the function is smooth at $(1 : 0 : . . . : 0)$ (from the inverse
  function theorem), and if we write the variables of $U_0$
  with $u_i := \frac{x_i}{x_0}$, we get that the curve is:
  \[
    f(1,u_1,\dots,u_n) =
    \sum_{|\alpha| = d}
    c_{\alpha}
    \vec u^{\alpha} = 0,
  \]
  and since as mentioned we require $\frac{\partial f}{\partial {u_i}}(0) = 0$,
  we get that $c_{\beta_i}$ where $\beta_i$ is the multi-index with $1$
  in the $i$th index and $0$ everywhere else (except the $0$ index in which
  it is $d$), must be $0$ for all $1 \le i \le n$.
  In total, we get that
  \[
    c_{(d,0,\dots,0)} =
    \beta_1 = \cdots = \beta_{n} = 0.
  \]
  Thus, the subspace of the coefficient space is of condimension $n+1$.
\end{solution}

\begin{exercise}
  For the curve $E$ given by
  \[
    y^2 z = (x - \alpha_1 z) (x - \alpha_2 z) (x - \alpha_3 z),
  \]
  where $\alpha_i$ are distinct
  and $O = (0 : 1 : 0)$ compute
  explicitly the sum $A + B$ for $A = (x_1 : y_1 : 1)$ and $B = (x_2 : y_2 : 1)$
  (you can assume $x_1 \neq x_2$).
  Verify that the operation is associative.
  % Hint: write a parametric equation of the line $AB$,
  % restrict the equation of $E$ to it and
  % then use the Vieta’s formula for the sum of three roots of a cubic polynomial.
  % Don’t be shy about using software for algebraic manipulations.
\end{exercise}
\begin{solution}
  We work in the affine chart $z = 1$.
  The line through $A$ and $B$ has the equation:
  \[
    \ell \colon y = \lambda x + \nu,
    \quad \text{where }
    \lambda = \frac{y_2 - y_1}{x_2 - x_1},
    \quad \text{and } \nu = y_1 - \lambda x_1.
  \]
  Once we subtitute $y = \lambda x + \nu$ into the euqation of the curve
  we get:
  \[
    (\lambda x + \nu)^2 = (x - \alpha_1 z) (x - \alpha_2 z) (x - \alpha_3 z),
  \]
  which is a cubic equation in $x$, whose $3$ roots are exactly the $x$-coordinates
  of the intersections of $\ell$ with $E$ (for a sanity check, indeed
  there are $3$ intersections from B\'ezout's theorem).
  Two of these roots are $x_1$ and $x_2$,
  and if we denote the third root by $x_3$,
  by Vieta's formula we get that:
  \[
    x_1 + x_2 + x_3 = \alpha_1 + \alpha_2 + \alpha_3 + \lambda^2,
  \]
  or simply $x_3 = \alpha_1 + \alpha_2 + \alpha_3 + \lambda^2 - x_2 - x_1$.
  The group law states that $A + B$ is the third intersection
  of the line $OX$ with $E$ (the one that is not $O$ nor $X$),
  or in other words (since the cubic curve is symmetric with respect to the
  $x$-axis),
  the reflection across the $x$-axis of $(x_3 : y_3 : 1)$.
  Thus we get that $A + B = (x_3 : -y_3 : 1)$.
  We can verify that this group law is associative in modern mathmatical
  sofware (such as SageMath).
\end{solution}

\begin{exercise}
  Verify that $O$ is the identity of the group law on $E$.
  Verify that the inverse $B$ of $A \in E$
  is the third intersection point of the line $AO$ with $E$.
\end{exercise}
\begin{solution}
  As we saw, in general the group law with resepct to $O = (0 : 1 : 0)$
  in the chart $z = 1$ states that $A+B$ is the reflection
  of the third intersection point of the line $AB$ with $E$ with respect to the $x$-axis.
  If $A = A$ and $B = O$ then we get the reflection of the third intersection
  point of $AO$ with $E$, which is in itself the reflection of $A$ with resepct
  to $x$-axis, and since reflection is an involution
  (and by that I mean that it is the inverse of itself) we get that $A + O = A$.

  Suppose that $B$ is the third intersection point of the line $AO$ with $E$.
  As we saw (since the cubic curve is symmetric with respect to the $x$-axis),
  it is simply the reflection of $A$ with respect to the $x$-axis so
  $B = (x_A,-y_A)$.
  We then have that $A + B$ is the reflection of the third intersection point
  of the line $AB$ with $E$ with respect to the $x$-axis,
  or in other words, the reflection of $O$ with respect to the $x$-axis,
  but this is just $O$ itself in the projective space!
  This shows that $A + B = O$ which completes the proof.
\end{solution}

\newpage

\section{Cubic curves as complex Lie groups. Lattice description. Elliptic functions.}

\begin{exercise}
  Prove that for every lattice $\Lambda$ there exists a constant $c$ such
  that the number of $l \in L$ with $k \le |l| < k + 1$ is at most $c k$
  for all $k \geq 1$.
  % Hint: Argue that the union of fundamental parallelograms of L centered at
  % such l lies in a certain annulus and thus has smaller area.
\end{exercise}
\begin{solution}
  Let every lattice point $l \in \Lambda$ a copy $F_l$
  of the fundamental parallelogram centered at $l$
  which is simply $F_l := l + F$.
  Note that $\set{F_l}_{l \in \Lambda}$ partition the plane.
  Define:
  \[
    S_k :=
    \set{l \in \Lambda \mid k \le |l| < k + 1},
  \]
  and let $U_k := \bigcup_{l \in S_k} F_l$.
  Each $F_{l}$ has area $\area(F)$,
  and so $\area(U_k) = |S_k| \cdot \area(F)$.
  Let $R := \diam(F)/2$.
  Since this is the maximal distance from the center of a parallelogram
  to any other point in the parallelogram we get:
  \[
    U_k \subset
    \set{z \in \C \mid k - R \le |z| \le k + 1 + R} =: A_k.
  \]
  We can compute using a known formula for the area of an annulus
  and see that:
  \[
    \area(A_k) = \pi\del{(k + 1 + R)^2 - (k - R)^2} = \cdots = \O(k),
  \]
  as $k \to \infty$.
  This implies that $\area(A_k) \le C k$ for some constant $C$ depending
  only on the lattice $\Lambda$.
  Since $U_k \subset A_k$ we get that:
  \[
    |S_k| \cdot \area(F) = \area(U_k) \le \area(A_k) \le C k,
  \]
  and so it follows that
  \[
    |S_k| \le \frac{C}{\area(F)} k =: ck,
  \]
  for all $k \geq 1$ for a sufficienly large $C$,
  which gives us a sufficiently large $c := \frac{C}{\area(F)}$.
\end{solution}

\begin{exercise}
  Prove that for any even elliptic function $f$ and any $a \in \half \Lambda$
  the order of pole or zero of $f$ at $a$ is even.
  % Hint: Consider the Laurent power
  % series of f at a and use f(2a −z) = f(−z) = f(z).
\end{exercise}
\begin{solution}
  It suffices to show this claim for half-periods in the fundamental
  parallelograms.
  Let $\omega$ be a period of $f$ and $z \in \C$.
  Since $f$ is even we have:
  \[
    f(\omega/2 + z) =
    f(\omega/2 + z - \omega) =
    f(-\omega/2 + z) =
    f(\omega/2 - z).
  \]
  This holds for any $z \in \C$, and if we look at what this means
  on the level of the Laurent expansions of $f$ at $a := \omega/2$ we get:
  \[
    \sum_{n=-m}^{\infty} c_n (z)^n =
    \sum_{n=-m}^{\infty} c_n (-z)^n =
    \sum_{n=-m}^{\infty} c_n (-1)^{n}(z)^n,
  \]
  and thus by uniqueness of the Laurent series we must have $c_n = (-1)^{n} c_n$,
  and thus when $n$ is odd $c_n = 0$.
  Since only positive powers appear in the Laurent series expansion,
  we can conclude if $\omega/2$ were a pole (or zero, in which case $m$
  is negative) its order must be even.
\end{solution}

\begin{exercise}
  Prove that every elliptic function can be uniquely written
  as a sum of an even and an odd elliptic functions.
\end{exercise}
\begin{solution}
  First we recall that $\wp$ is even and $\wp'$ is odd.
  Recall that we can write $f$ as the sum of the even and odd
  functions:
  \[
    f_{\text{even}}(z) := \frac{f(z) + f(-z)}{2} \tand
    f_{\text{odd}}(z) := \frac{f(z) - f(-z)}{2}.
  \]
  Since $f$ is $\Lambda$-periodic (where $\Lambda$ is the lattice of the
  elliptic function), we have that for any $\omega \in \Lambda$ that:
  \[
    f_{\text{even}}(z+\omega) =
    \frac{f(z+\omega) + f(-(z+\omega))}{2} =
    \frac{f(z) + f(-z)}{2} =
    f_{\text{even}}(z),
  \]
  and similarly,
  \[
    f_{\text{odd}}(z+\omega) =
    \frac{f(z+\omega) - f(-(z+\omega))}{2} =
    \frac{f(z) - f(-z)}{2} =
    f_{\text{odd}}(z),
  \]
  so both functions of the decomposition are $\Lambda$-periodic,
  and therefore elliptic.

  Suppose that this decomposition is not unique,
  so that we can also write $f$ as the sum of some $g_{\text{even}}$
  and $g_{\text{odd}}$.
  Then we get that:
  \[
    f_{\text{even}} - g_{\text{even}} =
    g_{\text{odd}} - f_{\text{odd}}.
  \]
  Since the left-hand side is even, and the right-hand side is odd,
  we must have that both sides are the constant zero function.
  This implies that the decomposition is unique and completes the proof.
\end{solution}

%% LOST CODE

\newpage

\section{Genus of complex algebraic curves and Riemann--Hurwitz formula}

\begin{exercise}
  Prove that for a nonconstant holomorphic function $f$ with $f(0) = 0$
  the number of solutions to $f(0) = \epsilon$ is equal,
  for all small nonzero $\epsilon$,
  to the smallest $k$ such that the $k$th derivative of $f$ at $z = 0$
  is nonzero.
  Hint: Since $f'(z)$ is holomorphic, there is a neighborhood where $f'(z)$ has
  no other zeros than perhaps $z = 0$.
  Thus any zeros of $f(z) - \epsilon$ will be simple.
  Then you can count the their number in a circle $|z| < \alpha$ by
  \[
    \frac{1}{2 \pi i} \oint_{|z| = \alpha} \frac{f'(z)}{f(z) - \epsilon} \dz.
  \]
  It is continuous in $\epsilon$ for $\epsilon$ close to $0$,
  and is equal to
  $\frac{1}{2 \pi i} \oint_{|z| = \alpha} \frac{f'(z)}{f(z)} \dz$.
\end{exercise}
\begin{solution}
  Without loss of generality assume that $0$ is the only zero of $f$
  (otherwise, we repeat the process for all zeros seperately and sum
  the results).
  Since $f$ is holomorphic at $0$ with a zero at $z = 0$ of order $k$,
  there exists a holomorphic function $g$ with $g(0) \neq 0$ such that
  \[
    f(z) = z^k g(z),
  \]
  for all $|z| \le a$ for some sufficiently small $a > 0$.
  Since $g(0) \neq 0$, this means that there exists some $b > a > 0$
  such that $g(z) \neq 0$ for all $0 < |z| < b$, and thus $f(z) \neq 0$
  on that punctured disc $D_b'(0)$.
  Since $f'$ is also holomorphic on that disc, from the same argument,
  there also exists some punctured disc $D_c'(0)$ such that
  $f'(z) \neq 0$ for $z \in D'_c(0)$.
  Let $\epsilon \neq 0$, and suppose that $f(z) = \epsilon$
  for some $z \in B := D_b'(0) \cap D_c'(0)$.
  Since $z \neq 0$, it follows that $f'(z) \neq 0$,
  and therefore all zeros of $f(z) - \epsilon$ in $B$
  are simple.

  Now take $\alpha > 0$ and $\epsilon > 0$
  sufficiently small such that $f(z) \neq 0$ and $f(z) \neq 0$
  on $0 < |z| < a$ and all solutions of $f(z) = \epsilon$
  (for small $|\epsilon|$) lie inside $0 < |z| < a$.
  Denote the number of solutions to $f(z) = \epsilon$ by $N(\epsilon)$.
  We can now compute $N(\epsilon)$ by using the argument principle:
  \[
    N(\epsilon) =
    \frac{1}{2 \pi i} \oint_{|z| = a} \frac{f'(z)}{f(z) - \epsilon} \dz.
  \]
  Note that $N(\epsilon)$ must be an integer,
  but also, since the integrand is a continuous function,
  $N(\epsilon)$ is continuous, and therefore it must be constant.
  It follows that:
  \[
    N(\epsilon) = N(0) =
    \frac{1}{2 \pi i} \oint_{|z| = a} \frac{f'(z)}{f(z)} \dz = k,
  \]
  since $f(z)$ is holomorphic, nonzero on values in $D_a(0)$ other than
  $z = 0$, and has a zero of order $k$ at $z = 0$.
  This completes the proof.
\end{solution}

\begin{exercise}
  Check the claim of Remark $14.2$ for the map $f \colon \CP^1 \to \CP^1$
  that sends $(x_0 : x_1)$ to $(x_0^n : x_1^n)$ for some positive integer $n$.
\end{exercise}
\begin{solution}
  Since $\CP^1$ is the Riemann sphere we use the standard result
  that $\Mer(\CP^1) = \C(z)$.
  Since here $X = Y = \CP^1$,
  it is clear that $\Mer(X)$ and $\Mer(Y)$ are fields of transcendence
  degree $1$ over $\C$ (Since $\mathrm{t.r.}\deg_{\C}(\C(z)) = 1$).

  Choose the affine coordinate $z = \frac{x_0}{x_1}$ on $\CP^1$.
  In this coordinate, the map becomes $z \mapsto z^n$.
  Then the pullback map $f^* \colon \Mer(\CP^1) \to \Mer(\CP^1)$
  is given by $f^*(\varphi)(z) = \varphi(z^n)$.
  Since $\Mer(\CP^1) = \C(z)$ we get that:
  \[
    f^*(\C(z)) = \C(z^n) \subseteq \C(z) = \Mer(X).
  \]
  Hence the field extension $\Mer(Y)/\Mer(X)$ is
  of degree $[\Mer(Y) : \Mer(X)] = [\C(z^n) : \C(z)] = n$,
  which is indeed the degree of the map $z \mapsto z^n$.
  Since algebraic topology is not a prerequisite to this class let
  us compute the degree of the map $z \mapsto z^n$.
  We have that for a generic point $(a : 1) \in \CP^1$
  with $a \neq 0$ the equation $z^n = a$ has exactly $n$ distinct solutions
  in the respective chart in $\CP^1$ which shows that
  $\deg f = n = [\Mer(Y) : \Mer(X)]$ and completes the proof.
\end{solution}

\begin{exercise}
  Prove the last claim of the section.
  Hint: First argue that for $k \geq 5$ the right hand side is at least $\half$.
  Then for $k = 4$ the right hand is positive,
  so at least one of $n_i$ is $3$ or more,
  which gives at least $-2 + \half + \half + \half + \frac{2}{3} = \frac{1}{6}$.
  If $k \le 2$, the right hand side is negative, so we must have $k = 3$, etc.
\end{exercise}
\begin{solution}
  Let us denote:
  \[
    R_k := -2 + \sum_{i=1}^{k} \frac{n_i - 1}{n_i}.
  \]
  Since $R_k$ must be positive and $n_i$ are integers,
  and each of $n_i$ is at least $2$,
  the minimal value of $R_k$ for $k \geq 5$ is achieved when $n_i = 2$
  for all $1 \le i \le k$.
  We get that for $k \geq 5$ we have:
  \[
    R_k \geq -2 + \sum_{i=1}^{5} \frac{2 - 1}{2} = \half > \frac{1}{42}.
  \]
  For $k = 4$, we still need $R_k > 0$,
  but all $n_i$ are integers and at least $2$.
  Since $-2 + \sum_{i=1}^{4} \half = 0$,
  we conclude that at least one $n_i$ is $3$ and since:
  \[
    R_4 \geq -2 + \half + \half + \half + \frac{2}{3} = \frac{1}{6} > \frac{1}{42},
  \]
  we also get that $R_4$ is at least $1/6$ so the inequality is again satisfied.
  For $k \le 2$ we see that:
  \[
    R_k =
    -2 + \sum_{i=1}^{k} \frac{n_i - 1}{n_i} \le
    -2 + \sum_{i=1}^{2} \frac{n_i - 1}{n_i} \le
    -2 + \frac{(n_1 - 1)(n_2 - 1)}{n_1 n_2} \le -1 < 0,
  \]
  so for $R_k$ to break the inequality, $k$ must be $3$.
  However, it does not break the inequality, since:
  \[
    R_3 =
    -2 + \del{1 - \frac{1}{n_1}} +
         \del{1 - \frac{1}{n_2}} +
         \del{1 - \frac{1}{n_3}} =
    1 -
      \underbrace{\del{\frac{1}{n_1} + \frac{1}{n_2} + \frac{1}{n_3}}}_{= S_{(n_1,n_2,n_3)}}.
  \]
  Since $n_1,n_2,n_3 \geq 2$ are integers, and our goal is to find
  the maximum value of $S_k$ such that it is less than $1$.
  It is then left to check options by reverse lexicographical order (right dictionary order),
  on the set $\set{(n_1,n_2,n_3) \mid n_i \in \Z_{\geq 2}}/\sim$
  where $\sim$ is the equivalence relation of permutations.
  We get that:
  \begin{align*}
    S_{(2,2,2)} &= \half + \half + \half > 1, \\
    S_{(2,3,2)} &= \half + \frac{1}{3} + \half > 1, \\
    S_{(2,3,3)} &= \half + \frac{1}{3} + \frac{1}{3} > 1, \\
    S_{(2,3,4)} &= \half + \frac{1}{3} + \frac{1}{4} > 1, \\
    S_{(2,3,5)} &= \half + \frac{1}{3} + \frac{1}{5} > 1, \\
    S_{(2,3,6)} &= \half + \frac{1}{3} + \frac{1}{6} = 1, \\
    S_{(2,3,7)} &= \half + \frac{1}{3} + \frac{1}{7} < 1,
  \end{align*}
  and thus $S_{(2,3,7)} = \half + \frac{1}{3} + \frac{1}{7} = \frac{41}{42}$,
  is our desired value, and $R_3 \geq \frac{1}{42}$ as wanted,
  which completes the proof.
\end{solution}

\newpage

\section{Blowup of a point in \texorpdfstring{$\C^2$}{C}.
  Birational equivalence of algebraic varieties.}

\begin{exercise}
  Check that the blowup $S$ of $\C^n$ at the origin can be identified
  with the line bundle $\O(-1) \to \CP^{n-1}$.
\end{exercise}
\begin{solution}
  Recall that $S$ is simply:
  \[
    S =
    \set{(z, [l]) \in \C^n \times \CP^{n-1} \mid z \in L},
  \]
  where $L$ is the line in $\C^n$ represented by the point $[l] \in \CP^{n-1}$.
  The tautological line bundle over $\CP^{n-1}$ is defined such that
  the fiber over a point $[l] \in \CP^{n-1}$ is the line $L \subset \C^n$
  that the point $[l]$ represents.
  In other words, a total space $E$ defined by:
  \[
    E = \set{([l],v) \in \CP^{n-1} \times \C^n \mid v \in L},
  \]
  where $L$ is the line in $\C^n$ represented by the point $[l] \in \CP^{n-1}$.
  The total space of the bundle also comes with the projection
  map $\O(-1) \to \CP^{n-1}$ given by $\pi([l],v) = [l]$.
  The identification is clearly given by the map $\psi \colon S \to E$
  defined by $\psi(z,[l]) = ([l],z)$.
  It is left to show that $\psi$ is an isomorphism of complex manifolds.
  Clearly $\psi$ is well defined, and it is also clearly a bijection
  as a reordering of the components of the product.
  Since both $\psi$ and $\psi^{-1}$ are a reordering of components of the products,
  and a reordering of components it is just the identity in local coordinates,
  both of these maps are in particular holomorphic.
  It follows that $\psi$ is a biholomorphism, which completes the proof.
\end{solution}

\begin{exercise}
  Let $W$ be the subset of $\CP^2 \times \CP^2$ given by the conditions
  that the points $(x_0 : x_1 : x_2) \times (y_0 : y_1 : y_2)$
  satisfy $x_0 y_0 = x_1 y_1 = x_2 y_2$.
  Prove that $W$ is a smooth surface.
  Verify that the projection maps $W \to \CP^2$ are birational
  and find the image of the exceptional locus (i.e. points
  in $\CP^2$ with more than one preimage point).
  \textbf{Remark:} $W$ is the closure of the graph of the Cremona transformation
  $(x_0 : x_1 : x_2) \mapsto (\frac{1}{x_0} : \frac{1}{x_1} : \frac{1}{x_2})$.
\end{exercise}
\begin{solution}
  The space $W$ is defined by the equations $x_0 y_0 = x_1 y_1$ and
  $x_1 y_1 = x_2 y_2$.
  In order to check smoothness, we verify that the Jacobian of the map
  on any affine chart is of full rank $2$.
  Consider the chart $U = \set{x_0 \neq 0 \stand y_1 \neq 0}$.
  Once we set $x_0 = y_1 = 1$, the equations become
  $y_0 = x_1$ and $x_1 = x_2 y_2$.
  Then $W$ is given by the zeros of the function:
  \[
    F(x_1,x_2,y_0,y_2) := (y_0 - x_1, x_1 - x_2 y_2).
  \]
  The Jacobian matrix is:
  \[
    \Dif F(x_1,x_2,y_0,y_2) =
    \begin{pmatrix}
      \frac{\partial F_1}{\partial x_1} &
      \frac{\partial F_1}{\partial x_2} &
      \frac{\partial F_1}{\partial y_0} &
      \frac{\partial F_1}{\partial y_2} \\
      \frac{\partial F_2}{\partial x_1} &
      \frac{\partial F_2}{\partial x_2} &
      \frac{\partial F_2}{\partial y_0} &
      \frac{\partial F_2}{\partial y_2}
    \end{pmatrix} =
    \begin{pmatrix}
      -1 & 0 & 1 & 0 \\
      -1 & -y_2 & 0 & -x_2
    \end{pmatrix},
  \]
  where $F_1$, $F_2$ are the first and second coordinates of the function $F$.
  The Jacobian is clearly of degree $2$ at every point,
  and thus $W$ is a smooth manifold of dimension $4 - 2 = 2$.

  Let $\pi_1 \colon W \to \CP^2$ be the projection $(x,y) \mapsto x$.
  On the quasi-projective (dense open set) variety $U = \set{[x_0 : x_1 : x_2] \in \CP^2 \mid
  x_0 x_1 x_2 \neq 0}$,
  the equations $x_0 y_0 = x_1 y_1 = x_2 y_2 = \lambda$ imply that $\lambda \neq 0$,
  thus $y_i = \frac{\lambda}{x_i}$, which means that $y$ is completely
  determined by $x$ in the following explicit way:
  \[
    y = [\frac{1}{x_0} : \frac{1}{x_1} : \frac{1}{x_2}].
  \]
  Since $\pi_1$ has a well-defined inverse on $U$ which is an dense open
  set in $\CP^2$, $\pi_1$ is birational.
  By symmetry, $\pi_2 \colon (x,y) \mapsto y$ is also birational.

  The exceptional locus of $\pi_1$ (we only consider $\pi_1$ without loss of generality),
  is by definition the points $x \in \CP^2$ where $\pi_1^{-1}(x)$ is not a single
  point.
  We already saw that points without zero coordinates are not in that exceptional
  locus, so let us consider points with exactly one coordinate zero.
  Consider for example $x = [0 : 1 : 1]$.
  The equations are $0 \cdot y_0 = 1 \cdot y_1 = 1 \cdot y_2 = \lambda$.
  Since $\lambda = 0$, we must have $y_1 = y_2 = 0$,
  which forces $y = [1 : 0 : 0]$, which is a single point,
  and thus $x$ is not in the exceptional locus.
  Next consider a point with two zero coordinates such as $x = [1 : 0 : 0]$.
  The equations become $1 \cdot y_0 = 0 \cdot y_1 = 0 \cdot y_2 = \lambda$.
  This forces $y_0 = \lambda = 0$.
  However, $y_1$ and $y_2$ are free and can be anything.
  This implies that
  \[
    \pi_1^{-1}(x) =
    \set{[1 : 0 : 0]} \times \set{[y_0,y_1,y_2] \in \CP^2 \mid y_0 = 0} \cong \CP^1.
  \]
  Thus $[1 : 0 : 0]$ has more than one preimage,
  and by symmetry and the previous arguments we conclude that
  The points in $\CP^2$ with more than one preimage are exactly
  the three vertices:
  \[
    \set{[1 : 0 : 0], [0 : 1 : 0], [0 : 0 : 1]}.
  \]
\end{solution}

\begin{exercise}
  Let $x^2 - y^3 = 0$ be a singular curve on $\C^2$.
  Compute its proper preimage $(15.2)$ on the blowup $S$ of $\C^2$ at $(0,0)$
  and verify that it is a smooth curve.
  \textbf{Hint:} Work in the coordinate charts on $S$.
\end{exercise}
\begin{solution}
  The blowup $S$ of $\C^2$ at the origin is covered by two charts $U_1$
  and $U_2$, where $S \subset \C^2 \times \CP^1$ with homogeneous coordinates
  $[u : v]$ for the $\CP^1$ factor.
  The first chart $U_1$ is where $u \neq 0$
  so we set $u = 1$ and let $t = v/u$ so that the coordinates are $(x,t)$
  for $x \in \C$ and the blowup map is $\pi(x,t) = (x,xt)$,
  where we must have $y = xt$.
  Similarly, on $U_2$ we have $v \neq 0$ so we set $v = 1$ and let $s = u/v$.
  The coordinates are $(y,s)$ for $y \in \C$ and the blowup map is $\pi(y,s) = (ys,y)$,
  where we must have $x = ys$.

  We now need to compute:
  \[
    \overline{\pi^{-1}(\C \setminus (\C \cap \pi(\mathrm{Exc})))} =
    \overline{\pi^{-1}(\C \setminus \set{0,0})}.
  \]
  In the first chart $U_1$ where $x \neq 0$ and $y = xt$
  we get that:
  \[
    x^2 - (xt)^3 = 0 \implies
    x^2 (1 - xt^3) = 0.
  \]
  We have $x \neq 0$ in this chart so we divide by $x^2$
  to find that:
  \[
    1 - xt^3 = 0 \implies
    x = \frac{1}{t^3}.
  \]
  Thus the set $\pi^{-1}(\C \setminus \set{0,0})$ in this chart is
  $\set{(x,t) \mid x = \frac{1}{t^3} \stand x \neq 0}$.
  The closure of this includes the point it this set when $x \to 0$,
  but this means $t \to \infty$, however the ``point'' $(0,\infty)$
  is not in $U_1$ so this set is already a closed set.
  
  In the chart $U_2$ where $y \neq 0$ we have $x = sy$ and thus:
  \[
    (sy)^2 - y^3 = 0 \implies
    y^2 s^2 - y^3 = 0 \implies
    y^2 (s^2 - y) = 0,
  \]
  but $y \neq 0$ so we divide by $y^2$ and get $s^2 = y$.
  Thus the set $\pi^{-1}(\C \setminus \set{0,0})$ in this chart is
  $\set{(s,y) \mid y = s^2 \stand y \neq 0}$.
  The closure of this set is this union the point when we take $y \to 0$,
  which means $s \to 0$, and thus it also includes the point $(0,0)$.

  The global preimage is now simply the union of these sets:
  \[
    \set{(s,y) \mid y = s^2 \stand y \neq 0} \cup \set{(0,0)} \cup
    \set{(x,t) \mid x = \frac{1}{t^3} \stand x \neq 0}.
  \]

  We now have to verify that the preimage is a smooth curve.
  We check the derivative of the local defining equation in $U_2$.
  Let $g(s,y) = s^2 - y$, then the partial derivatives are:
  \[
    \frac{\partial g}{\partial s} = 2s \stand
    \frac{\partial g}{\partial y} = -1.
  \]
  A singularity point can only exist when both derivatives vanish,
  but since the derivative with respect to $y$ is the constant $1$,
  this is never the case.
  Similarly, for $h(x,t) = 1 - xt^3$ the partial derivatives are:
  \[
    \frac{\partial h}{\partial x} = -t^3 \stand
    \frac{\partial h}{\partial t} = -3xt^2.
  \]
  A singularity point can only exist when both derivatives vanish,
  but this can only happen when $t = 0$.
  However when $t = 0$ in our equation we have $h(x,0) = 1 - 0 = 0$,
  so this point is not on our curve.
  This completes the proof.
\end{solution}

\newpage

\section{Finite group actions on affine algebraic varieties.}

\begin{exercise}[Molien series]
  Let $V$ be a complex finite-dimensional representation of a finite group $G$.
  It induces a grading-preserving action of $G$ on $\Sym^*(V) \cong \C[x_1,\dots,x_n]$.
  Prove that the Hilbert series of the ring of invariants $\Sym^*(V)^G$ can be
  computed by:
  \[
    \sum_{d=0}^{\infty} \dim_{\C} \del{\Sym^d(V)}^G t^d =
    \frac{1}{|G|} \sum_{g \in G} \prod_{i=1}^{n} \frac{1}{1 - \lambda_{g,i} t} =
    \frac{1}{|G|} \sum_{g \in G} \det(1 - t g)^{-1},
  \]
  where $\lambda_{g,1},\dots,\lambda_{g,n}$ are eigenvalues of $g$ on $V$.
  Hint: For a finite-dimensional $G$-representation $W$ we have
  $\dim_{\C} W^G = \frac{1}{|G|} \sum_{g \in G} \tr(g,W)$.
  Then use the result that all finite order elements in $\GL(V)$
  are diagonilizable.
\end{exercise}
\begin{solution}
  We can use the hint in order to conclude that:
  \[
    \sum_{d=0}^{\infty} \dim_{\C} \del{\Sym^d(V)}^G t^d =
    \sum_{d=0}^{\infty} \del{\frac{1}{|G|} \sum_{g \in G} \tr(g,\Sym^d(V))} t^d.
  \]
  Since the inner sum is finite we can swap the order of summation
  and get:
  \[
    \sum_{d=0}^{\infty} \dim_{\C} \del{\Sym^d(V)}^G t^d =
    \frac{1}{|G|} \sum_{g \in G} \del{\sum_{d=0}^{\infty} \tr(g,\Sym^d(V)) \cdot t^d}.
  \]
  Since $G$ is finite, all elements in $G$ have finite order,
  and thus all of them are diagonalizable.
  Let the eigenvalues of $g$ acting on $V$
  be $\lambda_{g,1},\dots,\lambda_{g,n}$.
  By the previous remark, there exists a basis $\set{v_1,\dots,v_n}$ of $V$
  such that $g(v_i) = \lambda_{g,i} v_i$,
  and thus the induced action on the symmetric algebra $\Sym^d(V)$ has
  a basis consisting of monomials $v_1^{\alpha_1} \cdots v_n^{\alpha_n}$
  where $\sum a_i = d$. The action on these monomials is:
  \[
    g \cdot \del{v_1^{\alpha_1} \cdots v_n^{\alpha_n}} =
    \del{\lambda_{g,1}^{\alpha_1} \cdots \lambda_{g,n}^{\alpha_n}}
      v_1^{\alpha_1} \cdots v_n^{\alpha_n}.
  \]
  The trace $\tr(g,\Sym^d(V))$ is simply the sum of these eigenvalues
  over all possible combinations of $\alpha_i$ which is:
  \[
    \tr(g,\Sym^d(V)) =
    \sum_{\alpha_1 + \cdots + \alpha_n = d}
      \lambda_{g,1}^{\alpha_1} \cdots \lambda_{g,n}^{\alpha_n}.
  \]
  Now go back to our original equation for the Hilbert series,
  and consider the inner (infinite) sum of a fixed element $g \in G$:
  \[
    \sum_{d=0}^{\infty} \tr(g,\Sym^d(V)) \cdot t^d =
    \sum_{d=0}^{\infty} \del{\sum_{\alpha_1 + \cdots + \alpha_n = d}
      \lambda_{g,1}^{\alpha_1} \cdots \lambda_{g,n}^{\alpha_n}} \cdot t^d.
  \]
  This is exactly the expansion of a product of geometric series:
  \[
    \sum_{d=0}^{\infty} \tr(g,\Sym^d(V)) \cdot t^d =
    \prod_{i=1}^{n} \del{\sum_{\alpha_i = 0}^{\infty} (\lambda_{g,i} t)^{\alpha_i}} =
    \prod_{i=1}^{n} \frac{1}{1 - \lambda_{g,i} t}.
  \]
  Finally, we recall that for any linear operator $A$,
  its determinant is just the product of its eigenvalues
  $\det(I - tA) = \prod (1 - t \lambda_i)$.
  Thus we get that:
  \[
    \prod_{i=1}^{n} \frac{1}{1 - t \lambda_{g,i}} =
    \frac{1}{\prod_{i=1}^{n} 1 - t \lambda_{g,i}} =
    \frac{1}{\det(1 - tg)}.
  \]
  We can now conclude that:
  \[
    \sum_{d=0}^{\infty} \dim_{\C} \del{\Sym^d(V)}^G t^d =
    \frac{1}{|G|} \sum_{g \in G} \frac{1}{\det(1 - tg)},
  \]
  which concludes the proof.
\end{solution}

\begin{exercise}
  Let $M$ be a Hermitean form on $\C^n$ given by $\ip{v}{w} = \bar w^T M v$.
  Prove that an endomorphism:
  \[
    g \colon v \mapsto Av
  \]
  satisfies $\ip{Av}{Aw} = \ip{v}{w}$ for all $v$ and $w$ if and only if
  $\bar A^T M A = M$.
  Prove that if for some invertible matrix $S$ we have $M = \bar S^T S$,
  then $S A S^{-1}$ is a unitary matrix.
\end{exercise}
\begin{solution}
  In this exercise we let $^*$ denote conjugate and transpose of a map.
  We have that:
  \[
    \ip{Av}{Aw} =
    (Aw)^* M (Av) =
    (w^* A^*) M (Av) =
    w^* (A^* M A) v.
  \]
  The right side of the equality we need to prove
  is $\ip{v}{w} = w^* M v$.
  It is then obvious that the equality is satisfied if and only if $M = A^* M A$.

  Next, assume that $M = S^* S$ for some invertible matrix $S$.
  We must show that $U = S A S^{-1}$ is a unitary matrix.
  In other words, we need to show that $U^* U = I$.
  We compute and get that:
  \[
    U^* = (S A S^{-1})^*
    (S^{-1})^* A^* S^*,
  \]
  and thus using the assumption that $M = S^* S$ we get:
  \[
    U^* U =
    (S^{-1})^* A^* S^* S A S^{-1} =
    (S^{-1})^* A^* M A S^{-1},
  \]
  and by using the identity $A^* M A = M$ this is just:
  \[
    U^* U =
    (S^{-1})^* I S^{-1} =
    (S^{-1})^* S^{-1} = I,
  \]
  which completes the proof.
\end{solution}

\begin{exercise}
  Prove that the conjugation action of quaternions preserves the subspace
  $\R i + \R j + \R k$ and the norm on it.
\end{exercise}
\begin{solution}
  Denote $V := \R i + \R j + \R k$,
  and recall that each quaternion $q \in \mathbb H$ can be written as
  $q = a + v$ where $a \in \R$ and $v = bi + cj + dk \in V$.
  We want to show that for any $0 \neq q \in \mathbb H$ and any $v \in V$
  we have that $u' := q u q^{-1} \in V$.
  Recall that $v \in V$ if and only if $\bar v = -v$, this is what we will
  actually show.
  We have that:
  \[
    \overline{u'} =
    \overline{q u q^{-1}} =
    \overline{q^{-1}} \overline u \overline{q}.
  \]
  Since in quaternions $q^{-1} = \frac{\overline q}{\norm{q}^2}$,
  we get that $\overline{q^{-1}} = \frac{q}{\norm{q}^2}$,
  and since $u \in V$ we get $\overline u = -u$.
  We can substitute these into the equation and get:
  \[
    \overline{u'} =
    \frac{-q}{\norm{q}^2} u \overline{q} =
    - q u \frac{\overline{q}}{\norm{q}^2},
  \]
  and finally we use $q^{-1} = \frac{\overline q}{\norm{q}^2}$ again and conclude:
  \[
    \overline{u'} =
    - q u \frac{\overline{q}}{\norm{q}^2} =
    - q u q^{-1} = -u',
  \]
  so $u' \in V$ which shows that the
  action of quaternions preserves the subspace $V := \R i + \R j + \R k$.

  Next, in order to show that it also preserves the norm,
  we recall that the norm on the quaternions in multiplicative,
  and then for $u \in V$ we get:
  \[
    \norm{q u q^{-1}} =
    \norm{q} \cdot \norm{u} \cdot \norm{q^{-1}} =
    \norm{q} \cdot \norm{u} \cdot \frac{\norm{\overline{q}}}{\norm{q}^2} =
    \norm{q} \cdot \norm{u} \cdot \frac{1}{\norm{q}} = \norm{u},
  \]
  where we use the trivial fact that $\norm{\overline{q}} = \norm{q}$
  in the third inequality.
\end{solution}

\newpage

\section{Resolutions of \texorpdfstring{$A_n$}{An}
  and \texorpdfstring{$D_4$}{Dfour} singularities.}

\begin{exercise}
  \label{ex:subgroups-of-sl}
  Prove that every cyclic subgroup of $\SL(2,\C)$
  can be conjugated to the subgroup:
  \[
    G = \ll g \gg,
    g(x,y) = (\xi x, \xi^{-1} y),
  \]
  where $\xi = e^{2\pi i/n}$.
\end{exercise}
\begin{solution}
  Let $H = \ll M \gg$ be a cyclic subgroup of $\SL(2,\C)$ of order $n > 1$.
  Since $M^n = I$, its eigenvalues must be primitive $n$th roots of unity.
  Since $M \in \SL(2,\C)$, its determinant must be
  $\det M = \lambda_1 \lambda_2 = 1$,
  so $\lambda_2 = \lambda_1^{-1} =: \lambda^{-1}$,
  where $\lambda_1$ and $\lambda_2$ are $M$'s eigenvalues,
  and $\lambda$ is a more convenient notation we will use to describe.
  If $\lambda = \lambda^{-1}$, then we have $\lambda^2 = 1$,
  and thus $\lambda = \pm 1$.
  However, in both cases, since we consider the matrices up to scalar
  multiplication either $n = 1$, or $M$ is not diagonalizable,
  in which case it is conjugate to a Jordan block with a nonzero
  eigenvalue, so it is not nilpotent, which is also a contradiction,
  so we conclude that $\lambda \neq \lambda^{-1}$.

  This implies that there exists a matrix $P' \in \GL(2,\C)$
  such that $P'^{-1} M P' = \diag(\lambda,\lambda^{-1})$,
  however while $P' \in \GL(2,\C)$,
  we can actually make sure it is in $\SL(2,\C)$
  by setting $P := (\sqrt{\det P'})^{-1} P' \in \SL(2,\C)$,
  and now we still have $P^{-1} M P = \diag(\lambda,\lambda^{-1})$.
  Note that $g(x,y) = (\xi x, \xi^{-1} y)$ corresponds
  to the matrix $M_g := \diag(\xi,\xi^{-1})$ where $\xi = e^{2\pi i/n}$.
  We have shown that $M$ is conjugate to $\diag(\lambda,\lambda^{-1})$.
  However since $\lambda$ is a primitive root of unity,
  we get $M^k = M_g$ for some $k \geq 1$,
  and clearly $M = (M_g)^{k'}$ for some $k'$ coprime to $n$,
  thus both groups contain the generators of each other and thus
  subgroup $H$ is conjugate to the group $G$.
\end{solution}

\begin{exercise}
  Verify that the proper image of $uv - w^n$ in the chart with coordinates
  $u$, $\frac{v}{u}$, and $\frac{w}{u}$ is smooth.
\end{exercise}
\begin{solution}
  First let us consider the first chart with coordinates
  $(u,\alpha,\beta) = (u,\frac{v}{u},\frac{w}{u})$.
  Under this notation we have $v = u \alpha$ and $w = u \beta$.
  Subtituting this into the original equation we get:
  \[
    u (u \alpha) - (u \beta)^n =
    u^2 \alpha - u^n \beta^n = 0.
  \]
  In order to find the proper preimage,
  we factor out the exceptional divisor.
  The lowest power of $u$ is $u^2$ and thus:
  \[
    u^2 (\alpha - u^{n-2} \beta^n) =
    0.
  \]
  Thus, the equation for the proper preimage of this chart is:
  \begin{equation}
    \label{eq:seven-four}
    \tilde f(u,\alpha,\beta) = \alpha - u^{n-2} \beta^n = 0.
  \end{equation}
  We know that at any singularity point,
  all the derivatives must vanish together,
  so let us compute the derivatives:
  \[
    \Dif \tilde f_1 =
    \begin{pmatrix}
      \diffp{{\tilde f}}{u} \\
      \diffp{{\tilde f}}{\alpha} \\
      \diffp{{\tilde f}}{\beta}
    \end{pmatrix} =
    \begin{pmatrix}
      1 \\
      \text{some value} \\
      \text{some value}
    \end{pmatrix} =
    \begin{pmatrix}
      0 \\ 0 \\ 0
    \end{pmatrix}.
  \]
  From the first equation $1 = 0$ it is clear that no point
  satisfies this equality, and so we get that the partial derivatives
  cannot vanish together,
  so there cannot be any singularities in this chart.
  In other words, the proper preimage is smooth in this chart.
\end{solution}

\begin{exercise}
  Check that the $D_3$ case is the same as the $A_3$ case in two ways,
  by looking at the group action, and by making a change of variables in
  $u^2 + wv^2 + w^2 = 0$.
\end{exercise}
\begin{solution}
  First we consider the equation:
  \[
    u^2 + wv^2 + w^2 = 0.
  \]
  We complete the square with respect to $w$ and get:
  \[
    u^2 + \del{w + \frac{v^2}{2}}^2 - \frac{v^4}{4} = 0.
  \]
  Let $X = u$ and $Y = i(w + \frac{v^2}{2})$, and $Z = \frac{u}{\sqrt[4]{4}}$.
  The equation becomes:
  \[
    X^2 - Y^2 - Z^4 = 0 \iff
    (X - Y)(X + Y) - Z^4 = 0,
  \]
  so if we define the coordinate change $U = X - Y$,
  $V - X + Y$, and $W = Z$ we get that
  \[
    UV - W^4 = 0,
  \]
  which is the equation for the $A_3$ singularity.

  We know that $A_3$ corresponds to the cyclic group $C_4$,
  so since $\xi = e^{2\pi i /4} = i$, the generator for the group
  is $\diag(i,-i)$ and so the map is
  \[
    \begin{pmatrix}
      i & 0 \\
      0 & -i
    \end{pmatrix}
    \begin{pmatrix}
      x \\ y
    \end{pmatrix} =
    \begin{pmatrix}
      ix \\ -iy
    \end{pmatrix}.
  \]
  Thus the ring of invariants is generated
  by $w = xy$ and $u = x^4$ and $v = y^4$.

  The generator for the group corresponding to $D_3$
  is
  \[
    g =
    \begin{pmatrix}
      0 & i \\
      i & 0
    \end{pmatrix},
  \]
  and this corresponds to the action $(x,y) \mapsto (yi,xi)$.
  Clearly the order of the group generated by this map is $4$,
  which is the same as in the previous calculation
  and by \Cref{ex:subgroups-of-sl} we get that
  those two group actions are conjugate, which completes the proof.
\end{solution}

\newpage

\section{Resolutions of \texorpdfstring{$A_n$}{An} and \texorpdfstring{$D_4$}{D4}
singularities}

\begin{exercise}
  Prove that every involution of $\CP^1$ that sends $0$ to $\infty$
  can be written in the form $z \mapsto \frac{c}{z}$ for some nonzero
  constant $c$.
  Then argue that by scaling $z$ we can rewrite it as $z \mapsto \frac{1}{z}$.
\end{exercise}
\begin{solution}
  First recall that every automorphism of the Riemann sphere
  is simply a M\"obius transformation:
  \[
    f(z) = \frac{az + b}{cz + d},
  \]
  for $a,b,c,d \in \C$ where $ad - bc \neq 0$.
  It is determined by its value on $3$ points,
  so if we just apply the conditions $f(0) = \infty$
  and $f(\infty)  = 0$ then we get the equations:
  \begin{align*}
    0 = f(\infty) = \frac{a}{c} \implies a = 0 \\
    \infty = f(0) = \frac{b}{d} \implies d = 0.
  \end{align*}
  Thus, $f$ is of the form:
  \[
    f(z) = \frac{b}{cz} =
    \frac{c'}{z},
  \]
  for some constant $c' := \frac{b}{c} \in \C^{\times}$.
  Since the function is an involution we also have the condition:
  \[
    z = f(f(z)) = f(c'/z) = \frac{c'}{\frac{c'}{z}} = z,
  \]
  which holds for all such $c' \in \C^{\times}$.
  Thus, an involution which satisfies the required conditions
  is always of the form $z \mapsto \frac{c'}{z}$ for some constant
  $c' \in \C^{\times}$.
  Therefore if we scale the map by a factor of $(c')^{-1}$,
  that is, define the map $g(z) := f(c^{-1}z)$,
  we get a function $g \colon z \mapsto \frac{(c')^{-1}c'}{z} = \frac{1}{z}$,
  which also sends $0$ to $\infty$ and $\infty$ to $0$,
  which completes the proof.
\end{solution}

\begin{exercise}
  Prove that every order three automorphism of $\CP^1$ can be written as
  $z \mapsto e^{2\pi/3}z$ in some coordinates.
  Hint: Consider the corresponding element in $\PGL(2,\C)$ and its
  eigenvalues (up to common scaling).
\end{exercise}
\begin{solution}
  Every automorphism of $\CP^1$ is a M\"obius transformation:
  \[
    f(z) = \frac{az + b}{cz + d},
  \]
  for $a,b,c,d \in \C$ where $ad - bc \neq 0$.
  The corresponding element in $\PGL(2,\C)$,
  is simply the matrix:
  \[
    \begin{pmatrix}
      a & b \\
      c & d
    \end{pmatrix} \longleftrightarrow
    \frac{az + b}{cz + d} =: A \in \PGL(2,\C).
  \]
  Now the condition is that it is of order three,
  which means that $A^{3} = \lambda I$ for some $\lambda \in \C^{\times}$
  and $I \in \PGL(2,\C)$.
  This implies we can scale the matrix $A$,
  that is, redefine it as $A := \sqrt[3]{\lambda} A$ so that
  $A^{3} = I$.
  Thus, the eigenvalues of $A$
  must be cube roots of unity.
  That is, if we denote the eigenvalues of $A$ by $\mu_1$ and $\mu_2$,
  we have $\mu_1,\mu_2 \in \set{1,\omega,\omega^2}$,
  where $\omega = e^{2\pi/3}$.
  
  Assume, by way of contradiction, that $\mu_1 = \mu_2$.
  In this case, the Jordan normal form of $A$ must be
  either:
  \[
    \begin{pmatrix}
      \mu_1 & 1 \\
      0 & \mu_1
    \end{pmatrix} \tor
    \begin{pmatrix}
      \mu_1 & 0 \\
      0 & \mu_1
    \end{pmatrix},
  \]
  however, we know that $A^3 = I$,
  so the first option can be discarded (powers of such Jordan
  blocks are matrices with elements that are powers of $\mu_1$
  with a binomial coefficient, so its order is infinite, which is
  a contradiction to the finite order of $A$).

  Now we know that $A$ is a diagonalizble matrix with two distinct
  eigenvalues (actually the latter implies the former in this case),
  so we first assume
  that they are $\mu_1 = 1$ and $\mu_2 = \omega$.
  This means that there exists a diagonalizing basis in which
  the transformation becomes:
  \[
    \begin{pmatrix}
      \omega & 0 \\
      0 & 1
    \end{pmatrix} \longleftrightarrow
    \frac{\omega z}{1} \quad \text{so the map is} \quad
    z \mapsto \omega z,
  \]
  where we recall $\omega = e^{2\pi/3}$,
  so this completes the proof for the $(\mu_1, \mu_2) = (1,\omega)$ case.

  Note that diagonalizble elements in $\PGL(2,\C)$ are determined
  by the ratio of their eigenvalues, we we do not need to check
  the $(\mu_1, \mu_2) = (\omega,\omega^2)$ case,
  and since the $(1,\omega^2)$ case gives a map conjugate to the map $(1,\omega)$
  (since $\omega^2 = \bar \omega$ and $\bar 1 = 1$),
  and conjugate maps are the same map under a change of coordinates,
  the proof is complete.
\end{solution}

\begin{exercise}
  Check that the proper preimage of $u^2 + v^3 + w^5 = 0$ has no singular
  points in the coordinate charts with coordinates $u$, $\frac{v}{u}$,
  $\frac{w}{u}$ and $\frac{u}{v}$, $v$, $\frac{w}{v}$.
\end{exercise}
\begin{solution}
  First let us consider the first chart with coordinates
  $(\alpha,\beta,\gamma) = (u,\frac{v}{u},\frac{w}{u})$.
  Under this notation we have $v = u \alpha$ and $w = u \beta$.
  Subtituting this into the original equation we get:
  \[
    u^2 + (u \alpha)^3 + (u \beta)^5 =
    u^2 + u^3 \alpha^3 + u^5 \beta^5 =
    0.
  \]
  In order to find the proper preimage,
  we factor out the exceptional divisor.
  The lowest power of $u$ is $u^2$ and thus:
  \[
    u^2 (1 + u \beta + u^3 \beta^5) =
    0.
  \]
  Thus, the equation for the proper preimage of this chart is:
  \begin{equation}
    \label{eq:eight-six}
    \tilde f_1(u,\alpha,\beta) = 1 + u \alpha^3 + u^3 \beta^5 = 0.
  \end{equation}
  We know that at any singularity point,
  all the derivatives must vanish together,
  so let us compute the derivatives:
  \[
    \Dif \tilde f_1 =
    \begin{pmatrix}
      \diffp{{\tilde f_1}}{u} \\
      \diffp{{\tilde f_1}}{\alpha} \\
      \diffp{{\tilde f_1}}{\beta}
    \end{pmatrix} =
    \begin{pmatrix}
      \alpha^3 + 3 u^2 \beta^5 \\
      3 u \alpha^2 \\
      5 u^3 \beta^4
    \end{pmatrix} =
    \begin{pmatrix}
      0 \\ 0 \\ 0
    \end{pmatrix}.
  \]
  From the second equation it clearly follows that either $u$ or $\alpha$
  must be zero, however if $u$ is zero then from \Cref{eq:eight-six}
  we have $1 = 0$ so points of this form are not in the proper preimage,
  and on the other hand if $\alpha$ is zero then either $u$ is zero,
  in which case as we saw the point is not in the proper preimage,
  or $\beta$ is zero, in which case we also get $1 = 0$ so these
  points are not in the proper preimage either.
  
  Since there are no points in the proper preimage $\tilde f_1 = 0$
  where all the partial derivatives vanish,
  there cannot be any singularities in this chart.
  In other words, the proper preimage is smooth in this chart.

  Next consider the chart $(\gamma,v,\delta)$ for $\gamma = \frac{u}{v}$
  and $\delta = \frac{w}{v}$.
  Under these coordinates we get that:
  \[
    (v \gamma)^2 + v^3 + (v \delta)^5 = 
    v^2 \gamma^2 + v^3 + v^5 \delta^5 =
    0.
  \]
  We then factor out the exceptional divisor $v^2$ and we get:
  \[
    (v \gamma)^2 + v^3 + (v \delta)^5 = 
    v^2 (\gamma^2 + v + v^3 \delta^5) = 0.
  \]
  This shows that the equation of the proper preimage is given by:
  \[
    \tilde f_2(\gamma,v,\delta) =
    \gamma^2 + v + v^3 \delta^5 = 0.
  \]
  Then in order to check for singularities we again check if the points
  on which the Jacobian vanishes are in the proper preimage,
  and we see that:
  \[
    \Dif \tilde f_2 =
    \begin{pmatrix}
      \diffp{{\tilde f_2}}{\gamma} \\
      \diffp{{\tilde f_2}}{v} \\
      \diffp{{\tilde f_2}}{\delta}
    \end{pmatrix} =
    \begin{pmatrix}
      2\gamma \\
      1 + 3v^2 \delta^5 \\
      5 v^3 \delta^4
    \end{pmatrix} =
    \begin{pmatrix}
      0 \\ 0 \\ 0
    \end{pmatrix}.
  \]
  From the third equation either $v$ or $\delta$ must be zero.
  In case $v$ is zero, the second equation becomes $1 = 0$
  which is impossible, and in case $\delta$ is zero,
  the seoncd equation again becomes $1 = 0$.
  This shows immediately that there are no singular points in the proper
  preimage.
  In other words, the proper preimage is smooth in this chart as well,
  which completes the proof.
\end{solution}

\newpage

\section{Overview of intersection theory. Surfaces.}

\begin{exercise}
  \label{ex:nine-one}
  Use $(19.1)$ in order to check that on the blowup $X$ of a point
  in $\CP^2$ there holds $(L')^2 = 0$.
  What is the geometric meaning of this?
\end{exercise}
\begin{solution}
  Recall that we have $[L] = [L'] + [E]$,
  or with some abuse of notation $L = L' + E$.
  We then get that:
  \[
    (L')^2 =
    (L - E)^2 =
    L^2 - 2 L \cdot E + E^2,
  \]
  and since we know already from the computations in class
  that $L^2 = 1$ and $L \cdot E = 0$ and $E^2 = -1$,
  we conclude that:
  \[
    (L')^2 = 1 - 2 \cdot 0 + (-1) = 0,
  \]
  which completes the first part of the exercise.

  When we blow of the surface at a point we pretty much pull the surface
  apart, and in this way lines that meet at $p$ no longer
  intersect each other, and they are sitting ``side by side''.
  The fact that $(L')^2 = 0$ is the fact that its self intersection
  is $0$, which means that the curve can be moved to a position where it
  does not intersect with itself at all.
  This is possible exactly because of the ``pulling apart'' that we
  get as a result of the blowup.
  We can also think of $L'$ as a fiber of a fibartion,
  since it is pretty much that just from the definition.
\end{solution}

\begin{exercise}
  \label{ex:nine-two}
  Let $C$ be a smooth curve on a smooth projective surface $S$
  and let $X \to S$ be the blowup of $S$ at a point $p \in S$.
  Denote by $C'$ the proper preimage of $C$ in $X$.
  Prove that $(C')^2 = C^2 - 1$.
  This generalizes the example of \Cref{ex:nine-one}.
\end{exercise}
\begin{solution}
  We have assume $p \in C$ otherwise $(C')^2 = C^2$.
  If $\pi \colon X \to S$ is the blowup map,
  and if the curve passes through $p$ at multiplicity $m$,
  the total transform $\pi^* C$ in the the Picard group of $X$
  is given by:
  \[
    [\pi^* C] = [C'] + m [E],
  \]
  and since $C$ is smooth, actually we have that the multiplicity
  $m = 1$, so we can deduce that:
  \[
    C' = \pi^* C - E,
  \]
  and therefore:
  \[
    (C')^2 =
    (\pi^* C - E)^2 =
    (\pi^* C)^2 - 2 (\pi^* C \cdot E) + E^2.
  \]
  We know that the pullback preserves self-intersection so
  $(\pi^* C)^2 = C^2$,
  we know that $\pi^* C$ and $E$ have zero intersections because
  of the blowup so $\pi^* C \cdot E = 0$,
  and $E^2 = -1$ since by the same argument from class
  we always have that the self-intesection of the exceptional
  divisor is $-1$.
  Substituting these into the equation we get:
  \[
    (C')^2 =
    (\pi^* C)^2 - 2 (\pi^* C \cdot E) + E^2 =
    C^2 - 2 \cdot 0 + (-1) = C^2 - 1,
  \]
  which completes the proof.
\end{solution}

\begin{exercise}
  Let $X$ be a blowup of $\CP^2$ in $5$ points $p_1,\dots,p_5$
  in general position.
  Prove that the proper preimage of any line through two of these points
  has self-intersection $(-1)$.
  Prove that the proper preimage of the unique conic that passes through
  all five points also has self-intersection $(-1)$.
\end{exercise}
\begin{solution}
  Let $L$ be a line in $\CP^2$ passing through two of the blown-up
  points, for example $p_1$ and $p_2$.
  In the original plane we have $L^2 = 1$ for example from Bezout's
  theorem.
  Similarly to what we did in \Cref{ex:nine-two},
  we know that the relationship between the proper transform $L'$
  and the total transform $\pi^* L$ is given by:
  \[
    L' = \pi^* L - E_1 - E_2,
  \]
  where $E_1$ and $E_2$ are the exceptional divisors at $p_1$
  and $p_2$ respectively.
  We only subtract $E_1$ and $E_2$ since by assumption the line
  only passes through these points.
  We get that:
  \[
    (L')^2 =
    (\pi^* L - E_1 - E_2)^2 =
    (\pi^* L)^2 + E_1^2 + E_2^2 - 2(\pi^* L \cdot E_1) -2(\pi^* L \cdot E_2)
    + 2(E_1 \cdot E_2).
  \]
  Using again the fact that the self-intersection of an exceptional
  divisor is $-1$, that $E_i \cdot \pi^*L = 0$ for $i=1,2$,
  that $(\pi^* L)^2 = 1$ since it intersects with itself once,
  and that $E_1 \cdots E_2 = 0$ (the lines live above different points
  so they do not intersect),
  we conclude that:
  \[
    (L')^2 =
    1 + (-1) + (-1) - 0 - 0 + 0 = -1,
  \]
  which completes the first part of the proof.

  Let $C$ be the unique conic passing through $p_1,\dots,p_5$.
  By Bezout's we have $C^2 = 2 \cdot 2 = 4$.
  The blowup relation is now:
  \[
    C' = \pi^* C - E_1 - E_2 - E_3 - E_4 - E_5,
  \]
  where $E_i$ corresponds to the exceptional divisor at $p_i$ for
  each $1 \le i \le 5$.
  Since the total transform $\pi^* C$ does not intersect
  with excpetional divisors, and again exceptional
  divisors at different points live in different space and also
  do not intersect, we can compute that:
  \[
    (C')^2 = (\pi^* C)^2 + \sum_{i=1}^{5} E_i^2,
  \]
  where $E_i^2 = -1$, which does not cancel out since
  the intersection of the exception divisor with itself is $-1$
  (also known result).
  We conclude that:
  \[
    (C')^2 = 4 + 5(-1) = -1,
  \]
  utilizing the fact we mentioned earlier that $(\pi^* C)^2 = C^2 = 4$
  from Bezout.
  This completes the proof.
\end{solution}

\newpage

\section{Rational Surfaces. Del Pezzo surfaces.}

\begin{exercise}
  Prove that for any vector bundle $W \to X$ and a line
  bundle $L \to X$ we have:
  \[
    \P(W \otimes L) \cong \P W.
  \]
  Conclude that the projectivization of $\O(a) \oplus \O(b) \to \CP^1$
  is isomorphic to the Hirzebruch surface $H_{|a-b|}$.
\end{exercise}
\begin{solution}
  Let $W \to X$ be a vector bundle of rank $r + 1$. 
  Its projectivization is:
  \[
    \P(W) = \set{\text{lines in fibers of $W$}} \to X.
  \]
  Similarly, for a line bundle $L \to X$, we can form $W \otimes L$,
  which has the same rank since $L$ is a line bundle.
  We want to show that $\P(W) \cong \P(W \otimes L)$ over $X$. 

  Consider a point $x \in X$. 
  The fiber of the projectivization is:
  \[
    \P((W \otimes L)_x) = \P(W_x \otimes L_x).
  \]
  Note that $L_x$ is a $1$-dimensional vector space,
  so tensoring it with $W_x$ does not change the lines in $W_x$. 
  Formally, this means that there exists a canonical isomorphism of
  projective spaces:
  \[
    \P(W_x \otimes L_x) \cong \P(W_x),
  \]
  given by $[v \otimes \ell] \mapsto [v]$,
  which is well-defined because scaling $\ell$ does not change the line
  $[v \otimes \ell]$. 

  These fiberwise isomorphisms vary smoothly with $x \in X$,
  therefore they glue to a bundle isomorphism:
  \[
    \P(W \otimes L) \cong \P(W)
  \]
  over $X$, which proves the first part of the exercise. 

  We now consider an application of this result with $X = \CP^1$
  and $W = \O(a) \oplus \O(b)$. 
  Without loss of generality we may assume $a \geq b$. 
  Then:
  \[
    W \cong (\O(a - b) \otimes \O(0)) \otimes \O(b),
  \]
  because 
  \[
    (\O(a - b) \otimes \O(0)) \otimes \O(b) \cong \O(a) \oplus \O(b).
  \]
  Now by applying the previous result we get:
  \[
    \P(\O(a) \oplus \O(b)) \cong 
    \P((\O(a - b) \otimes \O(0)) \otimes \O(b)) \cong 
    \P(\O(a - b) \otimes \O(0)).
  \]
  By definition, the Hirzebruch surface $H_n$ is:
  \[
    H_n := \P(\O(n) \oplus \O) \to \CP^1.
  \]
  Therefore:
  \[
    \P(\O(a) \oplus \O(b)) \cong \P(\O(a - b) \oplus \O) = H_{|a-b|}.
  \]
  The absolute value appears because in Hirzebruch surfaces we actually
  have $H_n \cong H_{-n}$ which is also why we can assume $a \geq b$. 
  This completes the proof.
\end{solution}

\begin{exercise}
  Show that the complement $H_n \setminus F_0$ is isomorphic to $\CP^1 \times C$.
  Use it to prove that $H_n \setminus (S_0 \cup F_0) \cong \C^2$.
\end{exercise}
\begin{solution}
  % A point in $H_n$ is a line $[v : w] \subseteq \O(n) \oplus \O$ over some
  % $x \in \CP^1$, and the fiber $F_0$ is all lines over $x = 0$. 
  %
  In order to describe $H_n \setminus F_0$ explicitly,
  let us consider the trivialization over $U := \CP^1 \setminus \set{0}$.
  On $U \cong \C$ the line bundle $\O(n)$ is trivial,
  so over $U$ we have:
  \[
    \O(n) \oplus \O \cong \C \oplus \C \cong \C^2.
  \]
  The projectivization of a trivial rank $2$ bundle is:
  \[
    \P(\C^2 \times U) \cong \CP^1 \times U,
  \]
  but $U = \CP^1 \setminus \set{0}$, so the fiber over $0$ is exactly removed. 
  Hence it follows that:
  \[
    H_n \setminus F_0 \cong 
    \CP^1 \times (\CP^1 \setminus \set{0}) \cong 
    \CP^1 \times \C,
  \]
  which proves the first part of the exercise. 

  For the next part, recall that $H_n \setminus F_0 \cong \CP^1 \times \C$,
  and that the section $S_0$ is given fiberwise by the line:
  \[
    [v : 0] \in \P(\C \oplus \C),
  \]
  which is a constant section of the trivial $\CP^1$ bundle. 
  Thus, under the previous
  identification $H_n \setminus F_0 \cong \CP^1 \otimes \C$,
  the section $S_0 \setminus F_0$ becomes:
  \[
    \{\text{a fixed point in } \CP^1\} \times \C.
  \]
  After choosing coordinates on the fiber, we may assume this fixed
  point is $\infty \in \CP^1$. Therefore,
  \[
    H_n \setminus (S_0 \cup F_0) =
    (H_n \setminus F_0) \setminus S_0 \cong 
    (\CP^1 \times \C) \setminus (\set{\infty} \times \C) \cong
    \C \times \C,
  \]
  hence $H_n \setminus (S_0 \cup F_0) \cong \C^2$, which completes the proof.
\end{solution}

\begin{exercise}
  A $(-1)$-curve in a smooth surface $X$ is by definition isomorphic to $\CP^1$
  and satisfies $E^2 = -1$.
  Use the adjunction formula to prove that it satisfies $K \cdot E = -1$
  where $K$ is the Weil divisor that corresponds to the canonical line
  bundle $K_X \to X$.
\end{exercise}
\begin{solution}
  For a smooth curve $E \subset X$, the adjunction formula states that:
  \[
    (K_X \otimes \O_X(E)) \vert_E \cong K_E,
  \]
  and when we take the degree of both sides,
  we have that the right hand side is $2g(E) - 2$ 
  where $g(E)$ is the genus of $E$,
  and the degree on the left
  hand side is $\det K_X \vert_E + \deg \O_X(E) \vert_E = (K + E) \cdot E$. 
  Thus,
  \[
    (K + E) \cdot E = 2g(E) - 2,
  \]
  Since $E \cong \CP^1$ we have $g(E) = 0$,
  therefore the formula reduces to:
  \[
    (K + E) \cdot E =
    K \cdot E + E^2 = -2,
  \]
  but we also have $E^2 = -1$.
  Hence we deduce that $K \cdot E = -1$.
\end{solution}

\newpage

\section{Kodaira dimension and classification of algebraic surfaces}

\begin{exercise}
  Count the parameters of rank $4$ lattices in $\C^2$,
  up to linear transformations of $\C^2$ to show that they form a
  four-dimensional family.
\end{exercise}
\begin{solution}
  Recall that a rank $4$ lattice $\Lambda \subset \C^2$ is a discrete subgroup
  \[
    \Lambda = \Z v_1 \oplus \cdots \oplus \Z v_4,
  \]
  such that $v_1,\dots,v_4$ span $\C^2 \cong \R^4$. 
  Thus, we can equivalently consider elements of $\GL(4,\R)$,
  which is a vector space of real dimension $16$. 
  Two bases define the same lattice if and only if they differ by
  an element of $\GL(4,\Z)$, however since this group is discrete,
  quotienting by it does not change the dimension. 
  That is,
  \[
    \{\text{rank-$4$ latrices in $\C^2$}\} \sim \GL(4,\R) / \GL(4,\Z)
  \]
  is still real $16$-dimensional. 
  We know need to mod out this space by linear transformations of $\C^2$,
  that is (by a result from class), elements of $\GL(2,\C)$. 
  This group has complex dimension $4$, but real dimension $8$,
  so after modding by it we are left with a space of real 
  dimension $16 - 8 = 8$, or simply complex dimension $4$.
  We may also notice that this is equivalent to classying the moduli space
  of complex tori since they are of the form $\C^2 / \Lambda$,
  so the moduli space is of complex dimension $4$. Hurray!
  A perfectly equivalent result that feels more meaningful.
\end{solution}

\begin{exercise}
  The weighted projective space $W\P(1,1,1,3)$ is defined as the quotient
  of $\C^4 \setminus \set{(0,0,0,0)}$ with coordinates $(x_0,x_1,x_2,t)$
  by the action of $\C^{\times}$ given by $(x_0,x_1,x_2,t) \mapsto
  (\lambda x_0, \lambda x_1, \lambda x_2, \lambda^3 t)$.
  As a set, $W\P(1,1,1,3)$ can be covered by the charts $U_i = \set{x_i \neq 0}$
  and $U_t = \set{t \neq 0}$ which gives it a structure of a (singular)
  algebraic variety.
  Prove that $U_{0,1,2}$ are in natural bijection to $\C^3$ and $U_t$
  is in bijection to the quotient of $\C^3$ by the group $\Z/3\Z$
  that acts by scaling coordinates by the third roots of $1$.
  Prove that the surface:
  \[
    t^2 = x_0^6 + x_1^6 + x_2^6
  \]
  is smooth and maps $2 : 1$ to $\CP^2$, with ramification divisor given by:
  \[
    0 = x_0^6 + x_1^6 + x_2^6.
  \]
\end{exercise}
\begin{solution}
  Consider the charts $U_i$ for $i=0,1,2$ and without loss of generality
  consider $U_0 = \set{x_0 \neq 0}$ specifically.
  We can use the $\C^{\times}$ action to normalize $x_0 = 1$
  by choosing $\lambda = x_0^{-1}$ and then:
  \[
    (x_0,x_1,x_2,t) \sim 
    \del{1, \frac{x_1}{x_0}, \frac{x_2}{x_0}, \frac{t}{x_0^3}}.
  \]
  Define the following coordinates:
  \[
    u_1 = \frac{x_1}{x_0} \tand
    u_2 = \frac{x_2}{x_0} \tand
    v = \frac{t}{x_0^3}.
  \]
  These are $3$ unconstrained complex numbers and thus $U_0 \cong \C^3$,
  and the same argument applies to all other charts $U_{1,2}$ as well.
  Hence $H_{0,1,2} \cong \C^3$ which completes the first part of the exercise. 

  Next consider the chart $U_t = \set{t \neq 0}$, in which in order to
  normalize $t = 1$ we must pick $\lambda$ such that $\lambda^3 = t^{-1}$.
  There are exactly three such choices of $\lambda$ which differ by 
  multiplication by a cube root of unity. 
  Note that after normalization:
  \[
    (x_0,x_1,x_2,t) \sim (\lambda x_0, \lambda x_1, \lambda x_2, 1),
  \]
  so by defining coordinates $y_i = \lambda x_i$, we see that changing
  $\lambda$ by a cube root of unity $\zeta$, has the effect of multiplying
  all $y_i$ by that $\zeta$. 
  Therefore, points are identified under:
  \[
    (y_0,y_1,y_2) \sim (\zeta y_0, \zeta y_1, \zeta y_2),
  \]
  where $\zeta$ is a cube root of unity. 
  Thus, $U_t \cong \C^3 / (\Z/3\Z)$,
  where $\Z/3\Z$ acts by scalar multiplication by a cube root of unity,
  which completes the second part of the exercise.

  Next we want to show that the surface $X$ given by:
  \[
    t^2 = x_0^6 + x_1^6 + x_2^6,
  \]
  is smooth. 
  On coordinates $U_{0,1,2}$ for example $U_0 \cong \C^3$,
  with coordinates $u_1,u_2,v$ as defined earlier we have:
  \[
    v^2 = 1 + u_1^6 + u_2^6.
  \]
  Define the function $F(u_1,u_2,v) = v^2 - (1 + u_1^6 + u_2^6)$. 
  We see that its partial derivatives are $2v$ and $-6u_1^5$ and $-6u_2^5$
  respectively. 
  A singular point would require that all derivatives vanish together
  so $v = u_1 = u_2 = 0$, but the point $(0,0,0)$ is not on the surface,
  so there are no singularities on the charts $U_{0,1,2}$. 

  In $U_t$ we have $t = 1$ so the equation becomes:
  \[
    1 = x_0^6 + x_1^6 + x_2^6.
  \]
  This is a smooth affine surface in $\C^3$ which is invariant under the 
  $\Z/3\Z$ action since the powers are multiples of $3$, and $\zeta^3 = 1$. 
  The smoothness follows from the same argument from earlier,
  it is clear that the only possible singulrity is at $(0,0,0)$ which is 
  not on the surface. Therefore $X$ is smooth in the last chart $U_t$ 
  so we conclude it is smooth. 

  Consider the map $\pi X \to \CP^2$ given by the projection:
  \[
    \pi \colon [x_0 : x_1 : x_2 : t] \mapsto [x_0 : x_1 : x_2].
  \]
  This is well-defined since the $\C^{\times}$ action scales $x_i$ 
  uniformly. 
  When 
  \[
    0 \neq t^2 = x_0^6 + x_1^6 + x_2^6,
  \]
  for a fiber $\pi^{-1}(a_0,a_1,a_2)$ we have exactly two solutions,
  $t = \pm \sqrt{a_0^6 + a_1^6 + a_2^6}$,
  so the degree of the map is $2$ as wanted.

  Ramification occurs where the map fails to be $2 : 1$,
  at the origin, so it forces the ramification divisor to be exactly
  \[
    0 = x_0^6 + x_1^6 + x_2^6,
  \]
  which completes the proof.
\end{solution}

\begin{exercise}
  One more accessible example of $\mathrm K3$ surfaces is that of surfaces $X$
  in $\CP^5$ cut out by three quadratic equations.
  Do the rough parameter count, taking into account that only the linear span
  of the three equations matters.
  Hint: Count the ordered bases of the space of generators,
  which overcounts the parameter space by $\dim \GL(3,\C) = 9$.
  If you know the formula for the dimension of Grassmannians,
  that could be used as well.
\end{exercise}
\begin{solution}
  A quadric in $\CP^5$ is a homogeneous degree-$2$ polynomial in $6$ variables. 
  The vector space of all quadrics has dimension:
  \[
    \binom{6 + 2 - 1}{2} =
    \binom{7}{2} = 21,
  \]
  so it is isomorphic to $\C^{21}$.
  A surface of the given type is cut out by a $3$-dimensional subspace
  $W \subset \C^{21}$, that is, a choice of three linearly independent 
  quadrics, but only their linear span matters.
  
  First, we overcount by counting ordered triplets of independent quadrics,
  which gives $3 \times 21 = 63$ parameters.
  As mentioned in the hint, this overcounts by $\dim \GL(3,\C) = 9$,
  so by subtracting we get $63 - 9 = 54$. 

  Two spaces are isomorphic if they differ by a projective linear change
  of coordinates of $\CP^5$. 
  By known results we have:
  \[
    \Aut(\CP^5) = \PGL(6,\C) \tand 
    \dim \PGL(6,\C) = 6^2 - 1 = 35.
  \]
  After modding by this we get a space of dimension $54 - 35 = 19$. 
  Thus, the family of $\mathrm K3$ surfaces obtained as complete intersections 
  of three quadrics in $\CP^5$ has $19$ complex parameters.
\end{solution}

\newpage

\section{Chern classes of vector bundles. Chern character. Euler sequence.}

\begin{exercise}
  Verify claim $(22.1)$.
\end{exercise}
\begin{solution}
  The formula we have is:
  \[
    c(\Sym^3(W)) = (1 + 3x_1)(1 + 2x_1 + x_2)(1 + x_1 + 2x_2)(1 + 3x_2),
  \]
  and we also found that $c_1(W) = x_1 + x_2 =: c_1$
  and $c_2(W) = x_1 x_2 =: c_2$.
  We notice that it may be smart to open the multiplication in pairs
  here:
  \[
    (1 + 3x_1)(1 + 3x_2) =
    1 + 3(x_1 + x_2) + 9 x_1 x_2 = 1 + 3c_1 + 9c_2,
  \]
  and
  \[
    (1 + 2x_1 + x_2)(1 + x_1 + 2x_2) =
    1 + 3(x_1 + x_2) + 2(x_1^2 + x_2^2) + 5 x_1 x_2.
  \]
  Since $x_1^2 + x_2^2 = c_1^2 - 2c_2$ we get that the product
  of this pair is simply:
  \[
    1 + 3c_1 + 2(c_1^2 - 2c_2) + 5c_2 = 1 + 3c_1 + 2c_1^2 + c_2,
  \]
  and now by a pure matter of computation and collecting the terms by
  their degree:
  \begin{align*}
    c(\Sym^3(W)) &=
    (1 + 3c_1 + 9c_2) (1 + 3c_1 + 2c_1^2 + c_2) \\ &=
    1 + 6c_1 + (11c_1^2 + 10c_2) + (6c_1^3 + 30c_1 c_2) + (18c_1^2c_2 + 9c_2^2).
  \end{align*}
\end{solution}

\begin{exercise}
  Prove that the first Chern class $c_1(W)$ of a vector bundle $W$ is the first
  Chern class of its top exterior power $\Lambda^{\rk W} W$.
\end{exercise}
\begin{solution}
  Let $\rk W = r$.
  We need to show that $c_1(W) = c_1(\Lambda^{r} W)$.
  We can use the ``black magic'' from class called the splitting principle.
  It states that if the bundle is over a paracompact space $X$,
  then there exists a space $X'$ (called the flag bundle) associated to $W$,
  and a map $p \colon X' \to X$ such that the induced cohomology homomorphism
  $p^* \colon H^*(X) \to H^*(Y)$ is injective, and the pullback bundle
  breaks up as a direct sum of line bundles
  $W' := p^*W = L_1 \oplus \cdots \oplus L_r$.
  We now get that,
  \[
    c_1(W') = c_1(L_1) + \cdots + c_1(L_r).
  \]
  The top exterior derivative is isomorphic to:
  \[
    \Lambda^{r} W' \cong L_1 \otimes \cdots \otimes L_r,
  \]
  since the exterior power of a direct sum of line bundles is their tensor
  product.
  From properties of Chern classes, we have that:
  \[
    c_1(L_1 \otimes \cdots \otimes L_r) = c_1(L_1) + \cdots + c_1(L_r).
  \]
  Thus, we get that $c_1(W') = c_1(\Lambda^{r} W')$.
  By injectivity of the pullback map $p^*$ and naturality/functoriality
  of the Chern classes (meaning the pullback diagram commutes), we get that
  $c_1(W) = c_1(\Lambda^{r} W)$ which completes the proof.
\end{solution}

\begin{exercise}
  Prove that the coefficients of the Todd class of $X$ can be written
  as polynomials in the Chern classes of $TX$, without using the dimension
  of $X$, and verify $(22.2)$.
  Hint: by looking at:
  \[
    \ln c(X) = \sum_{i} \ln(1 + x_i),
  \]
  we can write the sums $\sum_{i=1}^{\dim_{\C} X} x_i^k$ as polynomials in $c_j$.
  Then consider $\ln \Td(X)$.
\end{exercise}
\begin{solution}
  If $x_i$ are the Chern roots of $TX$ then we have:
  \[
    \ln c(X) = \sum_{i=1}^{r} \ln(1 + x_i),
  \]
  and by recalling the Taylor expansion of $\log(1 + x)$ which is:
  \[
    \ln(1 + x) = x - \frac{x^2}{2} + \frac{x^3}{3} - \cdots,
  \]
  and therefore,
  \[
    \ln c(X) = \sum_{k \geq 1} \frac{(-1)^{k-1}}{k} \sum_{i} x_i^k.
  \]
  However, $\ln c(X)$ is a formal power series whose coefficients are polynomials
  in the Chern classes $c_i$ and thus each $\sum_{i} x_i^k$ is a polynomial
  in the Chern classes $c_i$ that is independent of $\dim X$.

  Note that when considering the Todd class and taking the logarithm on both
  sides we get that:
  \[
    \ln \Td(X) = \sum_{i} \ln \del{\frac{x_i}{1 - e^{-x_i}}},
  \]
  but we can also expand $\frac{x}{1 - e^{-x}}$ in the same way:
  \[
    \frac{x}{1 - e^{-x}} = \frac{x}{2} + \frac{x^2}{12} - \frac{x^4}{720}
    + \cdots,
  \]
  and conclude that:
  \[
    \ln \Td(X) = \half \sum_{i} x_i + \frac{1}{12} \sum_{i} x_i^2
      - \frac{1}{720} \sum_{i} x_i^4 + \dots,
  \]
  but each sum $\sum_{i} x_i^k$ is already known to be a polynomial in
  in the Chern classes of $TX$ without using $\dim X$,
  so the coefficients of the Todd class of $X$ can be written
  as polynomials in the Chern classes of $TX$ which completes
  the first part of the exercise.

  In order to show $(22.2)$ we can now use the expansion of $e^x$
  and substitute $\ln \Td X$ and get:
  \[
    \Td(X) = 1 + \ln \Td(X) + \half (\ln \Td X)^2 + \text{higher order terms},
  \]
  and now we just need to collect the terms by their degree.
  The linear term is $\half c_1$, and the quadratic term is:
  \[
    \frac{1}{12} (c_1^2 - 2c_2) + \half \del{\half c_1}^2 =
    \frac{1}{12} c_1^2 - \frac{1}{6} c_2 + \frac{1}{8} c_1^2 =
    \frac{c_1^2 + c_2}{12}.
  \]
  Thus, we conclude that:
  \[
    \Td(X) = 1 + \frac{c_1}{2} + \frac{c_1^2 + c_2}{12} + \cdots,
  \]
  which completes the proof for $(22.2)$,
  and the exercise in general.
\end{solution}

\newpage

\section{Cohomology of vector bundles and Hirzebruch--Riemann--Roch formula.}

\begin{exercise}
  Prove that Euler characteristic is additive on short exact sequences.
  Hint: use $(23.1)$ and the general statement about alternating sum of
  dimensions of vector spaces in a long exact sequence.
\end{exercise}
\begin{solution}
  Given a short exact:
  \[\begin{tikzcd}
    0 & {W_1} & {W_2} & {W_3} & 0
    \arrow[from=1-1, to=1-2]
    \arrow[from=1-2, to=1-3]
    \arrow[from=1-3, to=1-4]
    \arrow[from=1-4, to=1-5]
  \end{tikzcd}\]
  we know from $(23.1)$ that there exists a long exact sequence in
  cohomology:
  \[
    \begin{tikzcd}
      \cdots \ar[r] & H^i(X,W_1) \ar[r] & H^i(X,W_2) \ar[r] & H^i(X,W_3)
        \ar[out=0, in=180, looseness=2, overlay, lld] \\
      & H^{i+1}(X,W_1) \ar[r] & H^{i+1}(X,W_2) \ar[r] & H^{i+1}(X,W_3) \ar [r] & \cdots
    \end{tikzcd}
  \]
  We don't know much about these coherent sheaf cohomologies,
  but they are vector spaces, and so they must satisfy the alternating
  sum of dimensions formula of a long exact sequence,
  we actually apply additivity of dimension of a short exact sequence of vector
  spaces multiple times (but certainly a finite amount of times, since
  these cohomology too eventually vanishes), and get by collecting
  terms of the same bundles together that:
  \[
    \sum_{i=0}^{\dim_{\C} X} (-1)^i \dim H^i(X,W_1) -
    \sum_{i=0}^{\dim_{\C} X} (-1)^i \dim H^i(X,W_2) +
    \sum_{i=0}^{\dim_{\C} X} (-1)^i \dim H^i(X,W_3) = 0.
  \]
  However, simply by substituting the definition of the Euler charateristic we
  conclude that:
  \[
    \chi(X,W_1) - \chi(X,W_2) + \chi(X,W_3) = 0,
  \]
  and thus $\chi(X,W_1) + \chi(X,W_3) = \chi(X,W_2)$,
  which completes the proof.
\end{solution}

\begin{exercise}
  \label{ex:eleven-two}
  Use Kodaira vanishing theorem in order to show that for all $k > - n - 1$
  we have $H^{> 0}(\CP^n,\O(k)) = 0$.
  Hint: Use Exercise $3$ of section $9$.
\end{exercise}
\begin{solution}
  The result from the exercise mentioned in the hint
  is $\O(-n-1) \cong K_{\CP^n}$, where $K_{\CP^n}$ is the canonical line
  bundle on $\CP^n$.
  Kodaira's vanishing theorem states that if $X$ is a smooth
  projective complex variety then if $L \to X$ is an ample line bundle
  we have $H^{> 0}(X,L \otimes K_X) = 0$.
  In particular, for $X = \CP^n$ we get that
  $H^{> 0}(\CP^n,L \otimes \O(-n-1)) = 0$.
  Note that we can write:
  \[
    \O(k) =
    \O(k + n + 1) \otimes K_{\CP^n} =
    \O(k + n + 1) \otimes \O(-n-1).
  \]
  We know that $\O(1)$ on $\CP^n$ is ample,
  and therefore every tensor power of it and in particular
  $\O(k + n + 1)$ is ample on $\CP^n$,
  and by applying Kodaira's vanishing theorem with $L = \O(k + n + 1)$
  we get that
  \[
    H^{> 0}(\CP^n,L \otimes \O(-n-1)) = H^{> 0}(\CP^n,\O(k)) = 0,
  \]
  which completes the proof.
\end{solution}

\begin{exercise}
  One can show (see Exercise $2$ of Section $7$) that for $k \geq 0$
  the global sections of $\O(k)$ on $\CP^n$ are homogeneous polynomials of degree
  $k$ in the homogeneous coordinates on $\CP^n$.
  Together with \Cref{ex:eleven-two}, show that this is consistent with the above
  calculation of $\chi(\CP^n,\O(k))$.
\end{exercise}
\begin{solution}
  Using \Cref{ex:eleven-two} we get that $H^{> 0}(\CP^n,\O(k)) = 0$
  and thus we get that:
  \[
    \chi(\CP^n,\O(k)) =
    \dim H^0(\CP^n,\O(k)) =
    \dim \Gamma(\CP^n,\O(k)),
  \]
  which as mentioned in the exercise we have:
  \[
    \Gamma(\CP^n,\O(k)) = 
      \set{\text{homogeneous polynomials of degree $k$ in $n + 1$ variables}}.
  \]
  This reduces the exercise into a counting game.
  The dimension of the number of sections is the number
  of different monomials $x_0^{a_0} x_1^{a_1} \cdots x_n^{a_n}$
  with $a_0 + \cdots a_n = k$,
  which, by a standard combinatorial result is $\binom{k+n}{n}$,
  so we have that:
  \[
    \chi(\CP^n,\O(k)) =
    \dim \Gamma(\CP^n,\O(k)) =
    \binom{k + n}{n} =
    \frac{(k + 1)(k + 2) \cdots (k + n)}{n!},
  \]
  which is exactly consistent with the calculation we did in class.
\end{solution}







\end{document}

